% Options for packages loaded elsewhere
\PassOptionsToPackage{unicode}{hyperref}
\PassOptionsToPackage{hyphens}{url}
%
\documentclass[
  10pt,
  a4paper,
]{article}
\usepackage{amsmath,amssymb}
\usepackage{iftex}
\ifPDFTeX
  \usepackage[T1]{fontenc}
  \usepackage[utf8]{inputenc}
  \usepackage{textcomp} % provide euro and other symbols
\else % if luatex or xetex
  \usepackage{unicode-math} % this also loads fontspec
  \defaultfontfeatures{Scale=MatchLowercase}
  \defaultfontfeatures[\rmfamily]{Ligatures=TeX,Scale=1}
\fi
\usepackage{lmodern}
\ifPDFTeX\else
  % xetex/luatex font selection
  \setmainfont[]{DejaVu Sans}
  \setmonofont[]{DejaVu Sans Mono}
\fi
% Use upquote if available, for straight quotes in verbatim environments
\IfFileExists{upquote.sty}{\usepackage{upquote}}{}
\IfFileExists{microtype.sty}{% use microtype if available
  \usepackage[]{microtype}
  \UseMicrotypeSet[protrusion]{basicmath} % disable protrusion for tt fonts
}{}
\makeatletter
\@ifundefined{KOMAClassName}{% if non-KOMA class
  \IfFileExists{parskip.sty}{%
    \usepackage{parskip}
  }{% else
    \setlength{\parindent}{0pt}
    \setlength{\parskip}{6pt plus 2pt minus 1pt}}
}{% if KOMA class
  \KOMAoptions{parskip=half}}
\makeatother
\usepackage{xcolor}
\usepackage[margin=0.83in]{geometry}
\usepackage{color}
\usepackage{fancyvrb}
\newcommand{\VerbBar}{|}
\newcommand{\VERB}{\Verb[commandchars=\\\{\}]}
\DefineVerbatimEnvironment{Highlighting}{Verbatim}{commandchars=\\\{\}}
% Add ',fontsize=\small' for more characters per line
\usepackage{framed}
\definecolor{shadecolor}{RGB}{248,248,248}
\newenvironment{Shaded}{\begin{snugshade}}{\end{snugshade}}
\newcommand{\AlertTok}[1]{\textcolor[rgb]{0.94,0.16,0.16}{#1}}
\newcommand{\AnnotationTok}[1]{\textcolor[rgb]{0.56,0.35,0.01}{\textbf{\textit{#1}}}}
\newcommand{\AttributeTok}[1]{\textcolor[rgb]{0.13,0.29,0.53}{#1}}
\newcommand{\BaseNTok}[1]{\textcolor[rgb]{0.00,0.00,0.81}{#1}}
\newcommand{\BuiltInTok}[1]{#1}
\newcommand{\CharTok}[1]{\textcolor[rgb]{0.31,0.60,0.02}{#1}}
\newcommand{\CommentTok}[1]{\textcolor[rgb]{0.56,0.35,0.01}{\textit{#1}}}
\newcommand{\CommentVarTok}[1]{\textcolor[rgb]{0.56,0.35,0.01}{\textbf{\textit{#1}}}}
\newcommand{\ConstantTok}[1]{\textcolor[rgb]{0.56,0.35,0.01}{#1}}
\newcommand{\ControlFlowTok}[1]{\textcolor[rgb]{0.13,0.29,0.53}{\textbf{#1}}}
\newcommand{\DataTypeTok}[1]{\textcolor[rgb]{0.13,0.29,0.53}{#1}}
\newcommand{\DecValTok}[1]{\textcolor[rgb]{0.00,0.00,0.81}{#1}}
\newcommand{\DocumentationTok}[1]{\textcolor[rgb]{0.56,0.35,0.01}{\textbf{\textit{#1}}}}
\newcommand{\ErrorTok}[1]{\textcolor[rgb]{0.64,0.00,0.00}{\textbf{#1}}}
\newcommand{\ExtensionTok}[1]{#1}
\newcommand{\FloatTok}[1]{\textcolor[rgb]{0.00,0.00,0.81}{#1}}
\newcommand{\FunctionTok}[1]{\textcolor[rgb]{0.13,0.29,0.53}{\textbf{#1}}}
\newcommand{\ImportTok}[1]{#1}
\newcommand{\InformationTok}[1]{\textcolor[rgb]{0.56,0.35,0.01}{\textbf{\textit{#1}}}}
\newcommand{\KeywordTok}[1]{\textcolor[rgb]{0.13,0.29,0.53}{\textbf{#1}}}
\newcommand{\NormalTok}[1]{#1}
\newcommand{\OperatorTok}[1]{\textcolor[rgb]{0.81,0.36,0.00}{\textbf{#1}}}
\newcommand{\OtherTok}[1]{\textcolor[rgb]{0.56,0.35,0.01}{#1}}
\newcommand{\PreprocessorTok}[1]{\textcolor[rgb]{0.56,0.35,0.01}{\textit{#1}}}
\newcommand{\RegionMarkerTok}[1]{#1}
\newcommand{\SpecialCharTok}[1]{\textcolor[rgb]{0.81,0.36,0.00}{\textbf{#1}}}
\newcommand{\SpecialStringTok}[1]{\textcolor[rgb]{0.31,0.60,0.02}{#1}}
\newcommand{\StringTok}[1]{\textcolor[rgb]{0.31,0.60,0.02}{#1}}
\newcommand{\VariableTok}[1]{\textcolor[rgb]{0.00,0.00,0.00}{#1}}
\newcommand{\VerbatimStringTok}[1]{\textcolor[rgb]{0.31,0.60,0.02}{#1}}
\newcommand{\WarningTok}[1]{\textcolor[rgb]{0.56,0.35,0.01}{\textbf{\textit{#1}}}}
\usepackage{longtable,booktabs,array}
\usepackage{calc} % for calculating minipage widths
% Correct order of tables after \paragraph or \subparagraph
\usepackage{etoolbox}
\makeatletter
\patchcmd\longtable{\par}{\if@noskipsec\mbox{}\fi\par}{}{}
\makeatother
% Allow footnotes in longtable head/foot
\IfFileExists{footnotehyper.sty}{\usepackage{footnotehyper}}{\usepackage{footnote}}
\makesavenoteenv{longtable}
\setlength{\emergencystretch}{3em} % prevent overfull lines
\providecommand{\tightlist}{%
  \setlength{\itemsep}{0pt}\setlength{\parskip}{0pt}}
\setcounter{secnumdepth}{5}
\ifLuaTeX
\usepackage[bidi=basic]{babel}
\else
\usepackage[bidi=default]{babel}
\fi
\babelprovide[main,import]{british}
\ifPDFTeX
\else
\babelfont{rm}[]{DejaVu Sans}
\fi
% get rid of language-specific shorthands (see #6817):
\let\LanguageShortHands\languageshorthands
\def\languageshorthands#1{}
\usepackage{fancyhdr}
\usepackage{fvextra}
\usepackage{xcolor}
\usepackage{microtype}
\definecolor{soloblue}{HTML}{256897}
\definecolor{solodark}{HTML}{24374A}
\usepackage{titlesec}
\titleformat{\section}{\Large\bfseries\color{soloblue}}{\thesection}{0.8em}{}
\titleformat{\subsection}{\large\bfseries\color{solodark}}{\thesubsection}{0.7em}{}
\setlength{\headheight}{15pt}
\pagestyle{fancy}
\fancyhf{}
\fancyhead[L]{\small\textcolor{soloblue}{SoloPhysics | Technical Development Handbook}}
\fancyhead[R]{\small\thepage}
\renewcommand{\headrulewidth}{0.4pt}
\AtBeginDocument{\renewcommand{\arraystretch}{1.16}}
\ifLuaTeX
  \usepackage{selnolig}  % disable illegal ligatures
\fi
\usepackage{bookmark}
\IfFileExists{xurl.sty}{\usepackage{xurl}}{} % add URL line breaks if available
\urlstyle{same}
\hypersetup{
  pdftitle={SoloPhysics: From a Mass-Spring Prototype to a 3D GLAD Renderer},
  pdfauthor={SoloPhysics development notes},
  pdflang={en-GB},
  hidelinks,
  pdfcreator={LaTeX via pandoc}}

\title{SoloPhysics: From a Mass-Spring Prototype to a 3D GLAD Renderer}
\usepackage{etoolbox}
\makeatletter
\providecommand{\subtitle}[1]{% add subtitle to \maketitle
  \apptocmd{\@title}{\par {\large #1 \par}}{}{}
}
\makeatother
\subtitle{An Obsidian + LaTeX Development Handbook for a Python
Programmer Learning C++}
\author{SoloPhysics development notes}
\date{31 October 2026}

\begin{document}
\maketitle

{
\setcounter{tocdepth}{2}
\tableofcontents
}
\section{Project Goals and
Chronology}\label{project-goals-and-chronology}

\subsection{Motivation}\label{motivation}

SoloPhysics began on \textbf{28 October 2026} with a long-term objective
of a BeamNG.drive-style deformable vehicle simulator, built
independently to learn C++, mathematical modelling and engine
development. The important design decision was to build a
\textbf{headless physics core first}, verify it numerically, and add
visualisation only after the results made sense. Inspiration from
node--beam vehicle physics is supported by
\href{https://www.beamng.com/game/about/physics/}{BeamNG's official
explanation}; the present project is an original prototype, not a
recreation of the full BeamNG solver.

\subsection{Development log --- 28 to 31 October
2026}\label{development-log-28-to-31-october-2026}

\begin{longtable}[]{@{}
  >{\raggedright\arraybackslash}p{(\columnwidth - 4\tabcolsep) * \real{0.3333}}
  >{\raggedright\arraybackslash}p{(\columnwidth - 4\tabcolsep) * \real{0.3333}}
  >{\raggedright\arraybackslash}p{(\columnwidth - 4\tabcolsep) * \real{0.3333}}@{}}
\toprule\noalign{}
\begin{minipage}[b]{\linewidth}\raggedright
Date
\end{minipage} & \begin{minipage}[b]{\linewidth}\raggedright
Milestone
\end{minipage} & \begin{minipage}[b]{\linewidth}\raggedright
Evidence / outcome
\end{minipage} \\
\midrule\noalign{}
\endhead
\bottomrule\noalign{}
\endlastfoot
\textbf{28 Oct} & Select C++20, GLM, CMake, Ninja, Git; headless-first
node--beam design & Development discussion \\
\textbf{28--29 Oct} & Create \texttt{Node}, \texttt{Beam},
\texttt{PhysicsWorld} and gravity / spring forces / integration &
Reviewed code fragments and user tests \\
\textbf{29 Oct} & Single suspended mass experiment; vary
stiffness/damping and plasticity & Equilibrium near \textbf{1.9019 m}
verified \\
\textbf{30 Oct} & Four-node, six-beam square dropped onto ground &
Position and velocity output supplied \\
\textbf{30 Oct} & Replace small-contact restitution with resting
threshold & Micro-bounce changed to nearly zero velocity \\
\textbf{30 Oct} & Add 2D GLFW/OpenGL renderer & User confirmed window
worked \\
\textbf{30 Oct} & Add camera controls, pause/reset/gravity toggles,
mouse-force dragging & User confirmed drag worked \\
\textbf{30 Oct} & Enable yielding/fracture, highlight nodes and draw
drag spring & User confirmed interactive damage worked \\
\textbf{31 Oct} & Replace 2D square with 3D cube, OpenGL 3.3 shaders,
orbit camera & Source bundle v0.4 verified \\
\textbf{31 Oct} & Initial 3D GLEW loader failed on user's setup; migrate
to GLAD 2 & User confirmed GLAD version worked \\
\end{longtable}

\begin{quote}
\textbf{Note:} Version labels v0.1--v0.4 are teaching/milestone labels
used during our development conversations, not a claim about Git tags or
formal releases. Exact commit IDs for the SoloPhysics project were not
provided here.
\end{quote}

\subsection{Why the order mattered}\label{why-the-order-mattered}

\begin{Shaded}
\begin{Highlighting}[]
\NormalTok{flowchart TD}
\NormalTok{ A[Physics data model] {-}{-}\textgreater{} B[Force accumulation]}
\NormalTok{ B {-}{-}\textgreater{} C[Time integration]}
\NormalTok{ C {-}{-}\textgreater{} D[Hanging mass verification]}
\NormalTok{ D {-}{-}\textgreater{} E[Square + collision]}
\NormalTok{ E {-}{-}\textgreater{} F[Stable resting contacts]}
\NormalTok{ F {-}{-}\textgreater{} G[2D live renderer]}
\NormalTok{ G {-}{-}\textgreater{} H[Interactive input + damage]}
\NormalTok{ H {-}{-}\textgreater{} I[3D camera + cube]}
\NormalTok{ I {-}{-}\textgreater{} J[GLAD loader migration]}
\NormalTok{ J {-}{-}\textgreater{} K[Future vehicle chassis]}
\end{Highlighting}
\end{Shaded}

The sequence prevents rendering bugs from being confused with simulation
bugs. For example, the 2D CLI output already proved the nodes were
moving before any shaders were written. Likewise, terminal damage counts
can distinguish physics failure from incorrect beam colouring.

\subsection{Fundamental architectural
choices}\label{fundamental-architectural-choices}

\begin{itemize}
\tightlist
\item
  \textbf{Simulation is separate from graphics:} the renderer receives a
  \texttt{const\ PhysicsWorld\&}; it reads state but does not integrate
  equations.
\item
  \textbf{Real units:} positions in metres, masses in kilograms,
  acceleration in m/s², force in newtons; spring coefficient in N/m,
  damping in N·s/m.
\item
  \textbf{Deterministic-size steps:} \texttt{1/240\ s} in interactive
  v0.3--v0.4. This is not a claim of fully deterministic cross-platform
  behaviour.
\item
  \textbf{Explicit topology:} beams store \emph{indices}, not owning
  pointers to nodes, so the world can store nodes in a
  \texttt{std::vector}.
\item
  \textbf{Per-beam failure state:} \texttt{broken} prevents the spring
  from transmitting force; \texttt{plasticallyDeformed} records
  historical deformation for display.
\end{itemize}

\subsection{Roadmap decision log}\label{roadmap-decision-log}

\textbf{Initially chosen:} straightforward mass--spring--damper with
semi-implicit Euler. \textbf{Deferred:} FEM, XPBD, constrained contacts,
self-collision, spatial acceleration structures, full vehicle dynamics,
solver calibration and multithreading. Their absence is an intentional
scope boundary rather than evidence they are solved.

See 18 - Limits and Future Roadmap and 19 - Bibliography and
Attribution.

\newpage

\section{Environment, Project Layout and
Build}\label{environment-project-layout-and-build}

\subsection{Hardware and operating
environment}\label{hardware-and-operating-environment}

The project is developed on \textbf{SoloLinux}, an Arch-derived Linux
setup, using a C++20 compiler (\texttt{g++} in the recorded build logs),
GLFW, GLM, Mesa/OpenGL, CMake, Ninja and later GLAD 2. The user was
running a Hyprland/Wayland desktop when the GLEW loader problem was
investigated.

\begin{Shaded}
\begin{Highlighting}[]
\FunctionTok{sudo}\NormalTok{ pacman }\AttributeTok{{-}S} \AttributeTok{{-}{-}needed}\NormalTok{ git python python{-}jinja glfw glm mesa cmake ninja}
\BuiltInTok{cd}\NormalTok{ \textasciitilde{}/Projects/SoloPhysics}
\FunctionTok{cmake} \AttributeTok{{-}S}\NormalTok{ . }\AttributeTok{{-}B}\NormalTok{ build{-}glad }\AttributeTok{{-}G}\NormalTok{ Ninja}
\FunctionTok{cmake} \AttributeTok{{-}{-}build}\NormalTok{ build{-}glad}
\ExtensionTok{./build{-}glad/solo\_physics}
\end{Highlighting}
\end{Shaded}

\begin{itemize}
\tightlist
\item
  \texttt{sudo\ pacman\ -S\ -\/-needed}: installs only missing packages.
\item
  \texttt{cmake\ -S\ .}: source directory is here.
\item
  \texttt{-B\ build-glad}: binary/build directory is a separate
  location.
\item
  \texttt{-G\ Ninja}: choose Ninja as the build-system generator.
\item
  \texttt{cmake\ -\/-build}: execute that generator with dependency
  checking.
\item
  \texttt{\&\&}: run the next command \textbf{only if the previous
  command succeeded}; prevents launching a stale executable after a
  failed compile.
\end{itemize}

\subsection{File structure}\label{file-structure}

\begin{Shaded}
\begin{Highlighting}[]
\NormalTok{SoloPhysics/}
\NormalTok{├── CMakeLists.txt}
\NormalTok{├── include/solophysics/}
\NormalTok{│   ├── Vector3D.hpp}
\NormalTok{│   ├── Node.hpp}
\NormalTok{│   ├── Beam.hpp}
\NormalTok{│   ├── PhysicsWorld.hpp}
\NormalTok{│   └── Renderer.hpp}
\NormalTok{├── src/}
\NormalTok{│   ├── main.cpp}
\NormalTok{│   ├── physics/PhysicsWorld.cpp}
\NormalTok{│   └── rendering/Renderer.cpp}
\NormalTok{└── build{-}glad/       \# generated; do not edit by hand}
\end{Highlighting}
\end{Shaded}

\textbf{Headers} (\texttt{.hpp}) describe interfaces and shared types.
\textbf{Implementation files} (\texttt{.cpp}) define functions.
\texttt{main.cpp} assembles the demo and owns the application loop. A
larger engine could split tests into a separate target and move the
engine into a library.

\subsection{Compilation mental model (from Python to
C++)}\label{compilation-mental-model-from-python-to-c}

Python often imports a module at runtime. C++ typically uses a
\textbf{preprocessor}, \textbf{compiler}, \textbf{linker}, then a native
executable:

\begin{Shaded}
\begin{Highlighting}[]
\NormalTok{flowchart LR}
\NormalTok{ A[main.cpp] {-}{-}\textgreater{} C[C++ compiler]}
\NormalTok{ B[Renderer.hpp / PhysicsWorld.hpp] {-}{-}\textgreater{} C}
\NormalTok{ D[Renderer.cpp] {-}{-}\textgreater{} E[C++ compiler]}
\NormalTok{ F[PhysicsWorld.cpp] {-}{-}\textgreater{} G[C++ compiler]}
\NormalTok{ C {-}{-}\textgreater{} H[Object files]}
\NormalTok{ E {-}{-}\textgreater{} H}
\NormalTok{ G {-}{-}\textgreater{} H}
\NormalTok{ H {-}{-}\textgreater{} I[Linker]}
\NormalTok{ J[GLFW / GLAD] {-}{-}\textgreater{} I}
\NormalTok{ I {-}{-}\textgreater{} K[solo\_physics]}
\end{Highlighting}
\end{Shaded}

\texttt{\#include\ \textless{}solophysics/Renderer.hpp\textgreater{}}
effectively makes declarations available to the current translation
unit. It is \textbf{not} a runtime package import. A missing declaration
causes a compile error; a declaration without a matching compiled
definition can cause a linker error.

\subsection{Git safeguards}\label{git-safeguards}

\begin{Shaded}
\begin{Highlighting}[]
\FunctionTok{git}\NormalTok{ status }\AttributeTok{{-}{-}short}
\FunctionTok{git}\NormalTok{ add CMakeLists.txt include src}
\FunctionTok{git}\NormalTok{ commit }\AttributeTok{{-}m} \StringTok{"Document current working GLAD renderer"}
\FunctionTok{git}\NormalTok{ log }\AttributeTok{{-}{-}oneline} \AttributeTok{{-}{-}max{-}count}\OperatorTok{=}\NormalTok{10}
\end{Highlighting}
\end{Shaded}

These are example commands; the exact history and branch state depend on
your repository. Commit before replacing generated upgrade files. See 16
- Errors, Diagnostics and Lessons.

\newpage

\section{C++ for a Python Developer}\label{c-for-a-python-developer}

\subsection{The biggest differences}\label{the-biggest-differences}

C++ is statically typed, compiled ahead of execution, and exposes object
lifetime, copies, references, pointers, stack/heap allocation and
generic templates more directly than Python. The SoloPhysics code uses
\textbf{C++20}, but many constructs predate C++20.

\begin{longtable}[]{@{}
  >{\raggedright\arraybackslash}p{(\columnwidth - 4\tabcolsep) * \real{0.3333}}
  >{\raggedright\arraybackslash}p{(\columnwidth - 4\tabcolsep) * \real{0.3333}}
  >{\raggedright\arraybackslash}p{(\columnwidth - 4\tabcolsep) * \real{0.3333}}@{}}
\toprule\noalign{}
\begin{minipage}[b]{\linewidth}\raggedright
Concept
\end{minipage} & \begin{minipage}[b]{\linewidth}\raggedright
Python
\end{minipage} & \begin{minipage}[b]{\linewidth}\raggedright
SoloPhysics C++
\end{minipage} \\
\midrule\noalign{}
\endhead
\bottomrule\noalign{}
\endlastfoot
Variable & \texttt{mass\ =\ 1.0} & \texttt{float\ mass\ =\ 1.0f;} \\
Constant-like binding & \texttt{DT\ =\ 1\ /\ 240} (convention) &
\texttt{constexpr\ double\ dt\ =\ 1.0\ /\ 240.0;} \\
Container & \texttt{nodes\ =\ {[}{]}} &
\texttt{std::vector\textless{}Node\textgreater{}\ nodes;} \\
Fixed-size collection & \texttt{positions\ =\ (...)} &
\texttt{std::array\textless{}Vector3D,\ 8\textgreater{}\ positions;} \\
Optional result & \texttt{None} or index &
\texttt{std::optional\textless{}std::size\_t\textgreater{}} \\
Function & \texttt{def\ step(dt):} & \texttt{void\ step(float\ dt);} \\
Exception & \texttt{raise\ ValueError(...)} &
\texttt{throw\ std::runtime\_error(...);} \\
Object attribute & \texttt{node.position} & \texttt{node.position} or
\texttt{nodePtr-\textgreater{}position} \\
Module organisation & \texttt{physics\_world.py} &
\texttt{PhysicsWorld.hpp} + \texttt{PhysicsWorld.cpp} \\
\end{longtable}

\begin{quote}
\textbf{Note:} Python analogy limitations Python names usually refer to
objects; C++ values can be copied by assignment. \texttt{Node\ a\ =\ b;}
copies \texttt{b}. \texttt{Node\&\ a\ =\ b;} is an alias to \texttt{b}.
\texttt{Node*\ p\ =\ \&b;} is a pointer containing its address.
\end{quote}

\subsection{From source to symbols}\label{from-source-to-symbols}

\begin{Shaded}
\begin{Highlighting}[]
\PreprocessorTok{\#include }\ImportTok{\textless{}glm/geometric.hpp\textgreater{}}
\KeywordTok{namespace}\NormalTok{ solophysics }\OperatorTok{\{}
    \KeywordTok{using}\NormalTok{ Vector3D }\OperatorTok{=}\NormalTok{ glm}\OperatorTok{::}\NormalTok{vec3}\OperatorTok{;}
\OperatorTok{\}}
\end{Highlighting}
\end{Shaded}

\begin{itemize}
\tightlist
\item
  \texttt{\#include} tells the preprocessor to make declarations
  available.
\item
  \texttt{namespace\ solophysics\ \{\ ...\ \}} groups symbols, rather
  like \texttt{solophysics.} names in Python.
\item
  \texttt{::} is the \textbf{scope resolution operator}:
  \texttt{glm::vec3} refers to \texttt{vec3} in namespace \texttt{glm}.
\item
  \texttt{using\ Vector3D\ =\ glm::vec3;} defines a type alias,
  \textbf{not a new class}.
\item
  \texttt{;} terminates most statements. \texttt{\{\}} delimit blocks
  instead of Python indentation.
\end{itemize}

\subsection{Values, references, const and
pointers}\label{values-references-const-and-pointers}

\begin{Shaded}
\begin{Highlighting}[]
\DataTypeTok{void}\NormalTok{ inspect}\OperatorTok{(}\AttributeTok{const}\NormalTok{ PhysicsWorld}\OperatorTok{\&}\NormalTok{ world}\OperatorTok{)} \OperatorTok{\{}
    \AttributeTok{const}\NormalTok{ Node}\OperatorTok{\&}\NormalTok{ a }\OperatorTok{=}\NormalTok{ world}\OperatorTok{.}\NormalTok{nodes}\OperatorTok{[}\DecValTok{0}\OperatorTok{];}
    \CommentTok{// a.mass = 5.0f; // ERROR: read{-}only reference}
\OperatorTok{\}}

\DataTypeTok{void}\NormalTok{ accelerate}\OperatorTok{(}\NormalTok{Node}\OperatorTok{\&}\NormalTok{ node}\OperatorTok{,} \AttributeTok{const}\NormalTok{ Vector3D}\OperatorTok{\&}\NormalTok{ force}\OperatorTok{)} \OperatorTok{\{}
\NormalTok{    node}\OperatorTok{.}\NormalTok{force }\OperatorTok{+=}\NormalTok{ force}\OperatorTok{;} \CommentTok{// changes the original node}
\OperatorTok{\}}
\end{Highlighting}
\end{Shaded}

\texttt{PhysicsWorld\&\ world} means \textbf{reference} (alias, no
copy). \texttt{const\ PhysicsWorld\&} makes the accessed object
read-only from this function. \texttt{Node*\ node} means
\textbf{pointer}. \texttt{\&node} in an expression takes an address, and
\texttt{*ptr} dereferences a pointer. The \texttt{*} in
\texttt{std::optional} dereferences the contained value too, but
\texttt{optional} is not itself a pointer.

\begin{Shaded}
\begin{Highlighting}[]
\NormalTok{GLFWwindow}\OperatorTok{*}\NormalTok{ window }\OperatorTok{=} \KeywordTok{nullptr}\OperatorTok{;}
\NormalTok{glfwMakeContextCurrent}\OperatorTok{(}\NormalTok{window}\OperatorTok{);}
\end{Highlighting}
\end{Shaded}

\texttt{nullptr} is the null-pointer literal; GLFW manages a window
behind this pointer. \textbf{Never} dereference a null pointer. The
renderer destroys the GLFW window during cleanup.

\subsection{Aggregates and brace
initialisation}\label{aggregates-and-brace-initialisation}

\begin{Shaded}
\begin{Highlighting}[]
\NormalTok{Node topLeft}\OperatorTok{;}
\NormalTok{topLeft}\OperatorTok{.}\NormalTok{position }\OperatorTok{=} \OperatorTok{\{{-}}\FloatTok{0.5}\BuiltInTok{f}\OperatorTok{,} \FloatTok{3.0}\BuiltInTok{f}\OperatorTok{,} \FloatTok{0.0}\BuiltInTok{f}\OperatorTok{\};}
\NormalTok{topLeft}\OperatorTok{.}\NormalTok{mass }\OperatorTok{=} \FloatTok{1.0}\BuiltInTok{f}\OperatorTok{;}
\end{Highlighting}
\end{Shaded}

\texttt{3.0f} is a \texttt{float} literal, while \texttt{3.0} is usually
a \texttt{double}. \texttt{\{x,\ y,\ z\}} supplies elements/constructor
inputs to a GLM vector. \texttt{Vector3D\{0.0f\}} constructs a zero
vector (as used in the project).

\subsection{Iteration and index types}\label{iteration-and-index-types}

\begin{Shaded}
\begin{Highlighting}[]
\ControlFlowTok{for} \OperatorTok{(}\NormalTok{Beam}\OperatorTok{\&}\NormalTok{ beam }\OperatorTok{:}\NormalTok{ world}\OperatorTok{.}\NormalTok{beams}\OperatorTok{)} \OperatorTok{\{}
\NormalTok{    beam}\OperatorTok{.}\NormalTok{breakStrain }\OperatorTok{=} \FloatTok{0.30}\BuiltInTok{f}\OperatorTok{;}
\OperatorTok{\}}

\ControlFlowTok{for} \OperatorTok{(}\BuiltInTok{std::}\NormalTok{size\_t i }\OperatorTok{=} \DecValTok{0}\OperatorTok{;}\NormalTok{ i }\OperatorTok{\textless{}}\NormalTok{ world}\OperatorTok{.}\NormalTok{nodes}\OperatorTok{.}\NormalTok{size}\OperatorTok{();} \OperatorTok{++}\NormalTok{i}\OperatorTok{)} \OperatorTok{\{}
    \AttributeTok{const}\NormalTok{ Node}\OperatorTok{\&}\NormalTok{ node }\OperatorTok{=}\NormalTok{ world}\OperatorTok{.}\NormalTok{nodes}\OperatorTok{[}\NormalTok{i}\OperatorTok{];}
\OperatorTok{\}}
\end{Highlighting}
\end{Shaded}

The first is a \textbf{range-based \texttt{for}} (like Python's
\texttt{for\ beam\ in\ beams}); \texttt{\&} avoids copying each
\texttt{Beam} and permits mutation. The second is an indexed loop;
\texttt{std::size\_t} is an unsigned size/index type, \texttt{++i}
increments it, \texttt{world.nodes{[}i{]}} indexes the vector. The
\texttt{i} variable exists \textbf{only inside the loop}---the exact
scope error that appeared in the project's 2D mouse-picking code.

\subsection{Automatic type deduction, templates and structured
bindings}\label{automatic-type-deduction-templates-and-structured-bindings}

\begin{Shaded}
\begin{Highlighting}[]
\AttributeTok{const} \BuiltInTok{std::}\NormalTok{array}\OperatorTok{\textless{}}\BuiltInTok{std::}\NormalTok{pair}\OperatorTok{\textless{}}\BuiltInTok{std::}\NormalTok{size\_t}\OperatorTok{,} \BuiltInTok{std::}\NormalTok{size\_t}\OperatorTok{\textgreater{},} \DecValTok{28}\OperatorTok{\textgreater{}}\NormalTok{ connections}\OperatorTok{\{\{}
    \OperatorTok{\{}\DecValTok{0}\OperatorTok{,} \DecValTok{1}\OperatorTok{\},} \OperatorTok{\{}\DecValTok{1}\OperatorTok{,} \DecValTok{2}\OperatorTok{\}} \CommentTok{/* ... other entries ... */}
\OperatorTok{\}\};}

\ControlFlowTok{for} \OperatorTok{(}\AttributeTok{const} \KeywordTok{auto}\OperatorTok{\&} \OperatorTok{[}\NormalTok{a}\OperatorTok{,}\NormalTok{ b}\OperatorTok{]} \OperatorTok{:}\NormalTok{ connections}\OperatorTok{)} \OperatorTok{\{}
\NormalTok{    world}\OperatorTok{.}\NormalTok{addBeam}\OperatorTok{(}\NormalTok{a}\OperatorTok{,}\NormalTok{ b}\OperatorTok{,}\NormalTok{ stiffness}\OperatorTok{,}\NormalTok{ damping}\OperatorTok{);}
\OperatorTok{\}}
\end{Highlighting}
\end{Shaded}

\begin{itemize}
\tightlist
\item
  \texttt{std::array\textless{}T,\ N\textgreater{}}: fixed-size sequence
  with \texttt{N} elements of type \texttt{T}.
\item
  \texttt{std::pair\textless{}A,B\textgreater{}}: two-value aggregate,
  like a typed 2-tuple.
\item
  \texttt{auto}: compiler infers the type; unlike Python, the type is
  still static.
\item
  \texttt{const\ auto\&\ {[}a,\ b{]}}: structured binding, unpacking the
  pair \textbf{without copying the pair}.
\item
  \texttt{\textless{}...\textgreater{}}: template arguments, not a
  less-than comparison in this context.
\end{itemize}

\subsection{Lambdas and capture
syntax}\label{lambdas-and-capture-syntax}

\begin{Shaded}
\begin{Highlighting}[]
\KeywordTok{auto}\NormalTok{ createWorld }\OperatorTok{=} \OperatorTok{[]()} \OperatorTok{\{}
\NormalTok{    PhysicsWorld world}\OperatorTok{;}
    \ControlFlowTok{return}\NormalTok{ world}\OperatorTok{;}
\OperatorTok{\};}

\KeywordTok{auto}\NormalTok{ applyMouseForce }\OperatorTok{=} \OperatorTok{[\&](}\NormalTok{PhysicsWorld}\OperatorTok{\&}\NormalTok{ physics}\OperatorTok{)} \OperatorTok{\{}
    \ControlFlowTok{if} \OperatorTok{(!}\NormalTok{grabbedNode}\OperatorTok{)} \ControlFlowTok{return}\OperatorTok{;}
\NormalTok{    Node}\OperatorTok{\&}\NormalTok{ node }\OperatorTok{=}\NormalTok{ physics}\OperatorTok{.}\NormalTok{nodes}\OperatorTok{[*}\NormalTok{grabbedNode}\OperatorTok{];}
\NormalTok{    node}\OperatorTok{.}\NormalTok{force }\OperatorTok{+=}\NormalTok{ mouseStiffness }\OperatorTok{*} \OperatorTok{(}\NormalTok{mouseTarget }\OperatorTok{{-}}\NormalTok{ node}\OperatorTok{.}\NormalTok{position}\OperatorTok{)}
                \OperatorTok{{-}}\NormalTok{ mouseDamping }\OperatorTok{*}\NormalTok{ node}\OperatorTok{.}\NormalTok{velocity}\OperatorTok{;}
\OperatorTok{\};}
\end{Highlighting}
\end{Shaded}

A lambda is an anonymous function (Python \texttt{lambda} analogy, but
C++ lambdas may contain multiple statements). \texttt{{[}{]}} is the
capture list. Empty \texttt{{[}{]}}: use no outside local variables.
\texttt{{[}\&{]}}: capture used enclosing local variables by
\textbf{reference}. The function body begins after \texttt{\{}.
\texttt{return} indicates the result; a \texttt{void} lambda has no
returned value. The user's callback is stored in
\texttt{std::function\textless{}void(PhysicsWorld\&)\textgreater{}} and
invoked from the physics timestep. \textbf{Lifetime rule:} captured
references must remain valid for as long as the callback uses them.

\subsection{Optional values}\label{optional-values}

\begin{Shaded}
\begin{Highlighting}[]
\BuiltInTok{std::}\NormalTok{optional}\OperatorTok{\textless{}}\BuiltInTok{std::}\NormalTok{size\_t}\OperatorTok{\textgreater{}}\NormalTok{ grabbedNode}\OperatorTok{;}
\ControlFlowTok{if} \OperatorTok{(}\NormalTok{grabbedNode}\OperatorTok{)} \OperatorTok{\{}
\NormalTok{    Node}\OperatorTok{\&}\NormalTok{ node }\OperatorTok{=}\NormalTok{ world}\OperatorTok{.}\NormalTok{nodes}\OperatorTok{[*}\NormalTok{grabbedNode}\OperatorTok{];}
\OperatorTok{\}}
\NormalTok{grabbedNode}\OperatorTok{.}\NormalTok{reset}\OperatorTok{();}
\end{Highlighting}
\end{Shaded}

\texttt{std::optional\textless{}T\textgreater{}} is an explicit `value
or no value' wrapper, similar to \texttt{int\ \textbar{}\ None} in
Python typing. \texttt{if\ (grabbedNode)} checks \textbf{presence}, not
whether the index is nonzero; index \texttt{0} is valid.
\texttt{*grabbedNode} retrieves it only after checking presence.
\texttt{.reset()} returns to the empty state.

\subsection{Classes, public/private, constructors and
cleanup}\label{classes-publicprivate-constructors-and-cleanup}

\begin{Shaded}
\begin{Highlighting}[]
\KeywordTok{class}\NormalTok{ Renderer }\OperatorTok{\{}
\KeywordTok{public}\OperatorTok{:}
\NormalTok{    Renderer}\OperatorTok{(}\DataTypeTok{int}\NormalTok{ width}\OperatorTok{,} \DataTypeTok{int}\NormalTok{ height}\OperatorTok{);}
    \OperatorTok{\textasciitilde{}}\NormalTok{Renderer}\OperatorTok{();}
\NormalTok{    Renderer}\OperatorTok{(}\AttributeTok{const}\NormalTok{ Renderer}\OperatorTok{\&)} \OperatorTok{=} \KeywordTok{delete}\OperatorTok{;}
\KeywordTok{private}\OperatorTok{:}
\NormalTok{    GLFWwindow}\OperatorTok{*}\NormalTok{ window }\OperatorTok{=} \KeywordTok{nullptr}\OperatorTok{;}
\OperatorTok{\};}
\end{Highlighting}
\end{Shaded}

\begin{itemize}
\tightlist
\item
  \texttt{class} defaults to private members, \texttt{struct} defaults
  to public members.
\item
  \texttt{public:} exposes methods to callers; \texttt{private:} hides
  implementation details.
\item
  \texttt{Renderer(...)} is a \textbf{constructor}, called on object
  creation.
\item
  \texttt{\textasciitilde{}Renderer()} is a \textbf{destructor}, called
  during destruction.
\item
  \texttt{=\ delete} prohibits copying the renderer; two copies owning
  one GLFW window could double-free GPU/window resources.
\item
  RAII (\emph{Resource Acquisition Is Initialization}) ties cleanup to
  object lifetime. Python has analogous context managers
  (\texttt{with}), but destruction timing differs.
\end{itemize}

\subsection{Error handling and
ownership}\label{error-handling-and-ownership}

\begin{Shaded}
\begin{Highlighting}[]
\ControlFlowTok{try} \OperatorTok{\{}
\NormalTok{    Renderer renderer}\OperatorTok{(}\DecValTok{1200}\OperatorTok{,} \DecValTok{800}\OperatorTok{);}
    \CommentTok{// application loop}
\OperatorTok{\}} \ControlFlowTok{catch} \OperatorTok{(}\AttributeTok{const} \BuiltInTok{std::}\NormalTok{exception}\OperatorTok{\&}\NormalTok{ e}\OperatorTok{)} \OperatorTok{\{}
    \BuiltInTok{std::}\NormalTok{cerr }\OperatorTok{\textless{}\textless{}} \StringTok{"SoloPhysics error: "} \OperatorTok{\textless{}\textless{}}\NormalTok{ e}\OperatorTok{.}\NormalTok{what}\OperatorTok{()} \OperatorTok{\textless{}\textless{}} \CharTok{\textquotesingle{}}\SpecialCharTok{\textbackslash{}n}\CharTok{\textquotesingle{}}\OperatorTok{;}
    \ControlFlowTok{return} \DecValTok{1}\OperatorTok{;}
\OperatorTok{\}}
\end{Highlighting}
\end{Shaded}

The \texttt{catch} catches standard-library exception types derived from
\texttt{std::exception}. \texttt{std::cerr} is the error output stream;
\texttt{\textless{}\textless{}} inserts text into the stream. Objects
with automatic lifetime clean up as the exception propagates.
Constructors must themselves clean up resources they acquired if
construction throws before completion.

\subsection{Syntax you will meet in the
renderer}\label{syntax-you-will-meet-in-the-renderer}

\begin{longtable}[]{@{}
  >{\raggedright\arraybackslash}p{(\columnwidth - 2\tabcolsep) * \real{0.5000}}
  >{\raggedright\arraybackslash}p{(\columnwidth - 2\tabcolsep) * \real{0.5000}}@{}}
\toprule\noalign{}
\begin{minipage}[b]{\linewidth}\raggedright
Syntax
\end{minipage} & \begin{minipage}[b]{\linewidth}\raggedright
Meaning
\end{minipage} \\
\midrule\noalign{}
\endhead
\bottomrule\noalign{}
\endlastfoot
\texttt{GLenum} / \texttt{GLuint} & OpenGL API-defined integer-like
types \\
\texttt{sizeof(Vertex)} & Byte size of a vertex structure \\
\texttt{offsetof(Vertex,\ colour)} & Byte offset of member
\texttt{colour} inside \texttt{Vertex} \\
\texttt{reinterpret\_cast\textless{}void*\textgreater{}(...)} & Explicit
low-level cast for the OpenGL vertex offset API \\
\texttt{R"GLSL(\ ...\ )GLSL"} & Raw string literal containing shader
source \\
\texttt{const\ char*} & Pointer to a null-terminated character
sequence \\
\texttt{static\_cast\textless{}float\textgreater{}(x)} & Explicit
checked-category conversion, e.g.~\texttt{double} to \texttt{float} \\
\texttt{std::clamp(x,\ low,\ high)} & Bound value within interval \\
\texttt{std::nullopt} & Empty optional value \\
\texttt{\#pragma\ once} & Widely supported header multiple-inclusion
protection \\
\texttt{constexpr} & Compile-time evaluable constant when applicable \\
\end{longtable}

\subsection{Pointer/reference traps for Python
developers}\label{pointerreference-traps-for-python-developers}

\begin{enumerate}
\def\labelenumi{\arabic{enumi}.}
\tightlist
\item
  A reference cannot be reseated after binding.
  \texttt{Node\&\ a\ =\ nodes{[}0{]}} will always refer to that node
  object during its lifetime.
\item
  \texttt{std::vector} can reallocate on growth, invalidating
  references/pointers to elements. This is why a beam stores indices,
  with explicit bounds checks.
\item
  \texttt{const} is enforced by the compiler in ordinary code; it isn't
  a runtime immutable proxy.
\item
  Ownership of C pointers to OpenGL/GLFW resources must be managed
  explicitly; don't assume Python-style garbage collection.
\item
  A \texttt{std::function} holding \texttt{{[}\&{]}} captures can create
  dangling references if its lifetime exceeds those local variables.
\end{enumerate}

See 20 - C++ Quick Reference for an exam-style sheet and 17 -
Verification Tests and Exercises for practice.

\newpage

\section{Core Data Model and Source
Architecture}\label{core-data-model-and-source-architecture}

\subsection{Core model}\label{core-model}

The simulation is a graph: \textbf{nodes are vertices}, and
\textbf{beams are undirected edges}. Each node has state (mass,
position, velocity, accumulated force); each beam refers to two node
indices and describes its constitutive behaviour.

\begin{Shaded}
\begin{Highlighting}[]
\NormalTok{classDiagram}
\NormalTok{class Node \{}
\NormalTok{  +Vector3D position}
\NormalTok{  +Vector3D velocity}
\NormalTok{  +Vector3D force}
\NormalTok{  +float mass}
\NormalTok{  +bool fixed}
\NormalTok{\}}
\NormalTok{class Beam \{}
\NormalTok{  +size\_t nodeA}
\NormalTok{  +size\_t nodeB}
\NormalTok{  +float restLength}
\NormalTok{  +float stiffness}
\NormalTok{  +float damping}
\NormalTok{  +float yieldStrain}
\NormalTok{  +float breakStrain}
\NormalTok{  +float plasticity}
\NormalTok{  +bool broken}
\NormalTok{  +bool plasticallyDeformed}
\NormalTok{\}}
\NormalTok{class PhysicsWorld \{}
\NormalTok{  +vector\textasciitilde{}Node\textasciitilde{} nodes}
\NormalTok{  +vector\textasciitilde{}Beam\textasciitilde{} beams}
\NormalTok{  +Vector3D gravity}
\NormalTok{  +step(dt)}
\NormalTok{  +addNode(node)}
\NormalTok{  +addBeam(a,b,k,c)}
\NormalTok{\}}
\NormalTok{class Renderer \{}
\NormalTok{  +render(world,...)}
\NormalTok{  +pickNode(world,radius)}
\NormalTok{  +updateCamera()}
\NormalTok{\}}
\NormalTok{PhysicsWorld "1" o{-}{-} "many" Node}
\NormalTok{PhysicsWorld "1" o{-}{-} "many" Beam}
\NormalTok{Renderer ..\textgreater{} PhysicsWorld : reads const reference}
\end{Highlighting}
\end{Shaded}

\subsection{\texorpdfstring{\texttt{Vector3D.hpp}: descriptive
alias}{Vector3D.hpp: descriptive alias}}\label{vector3d.hpp-descriptive-alias}

Recorded design:

\begin{Shaded}
\begin{Highlighting}[]
\PreprocessorTok{\#pragma once}
\PreprocessorTok{\#include }\ImportTok{\textless{}glm/vec3.hpp\textgreater{}}
\KeywordTok{namespace}\NormalTok{ solophysics }\OperatorTok{\{}
    \KeywordTok{using}\NormalTok{ Vector3D }\OperatorTok{=}\NormalTok{ glm}\OperatorTok{::}\NormalTok{vec3}\OperatorTok{;}
\OperatorTok{\}}
\end{Highlighting}
\end{Shaded}

\textbf{Interpretation:} \texttt{Vector3D} and \texttt{glm::vec3} denote
the same underlying type. GLM supplies vector addition, scalar
multiplication, dot products and lengths. Using an alias makes a future
math-library replacement easier conceptually, though replacing GLM would
require more work than just the alias.

\subsection{\texorpdfstring{\texttt{Node.hpp}: mass
particle}{Node.hpp: mass particle}}\label{node.hpp-mass-particle}

Recorded prototype (illustrative reconstruction of the earlier header):

\begin{Shaded}
\begin{Highlighting}[]
\KeywordTok{struct}\NormalTok{ Node }\OperatorTok{\{}
\NormalTok{    Vector3D position}\OperatorTok{\{}\FloatTok{0.0}\BuiltInTok{f}\OperatorTok{\};}
\NormalTok{    Vector3D velocity}\OperatorTok{\{}\FloatTok{0.0}\BuiltInTok{f}\OperatorTok{\};}
\NormalTok{    Vector3D force}\OperatorTok{\{}\FloatTok{0.0}\BuiltInTok{f}\OperatorTok{\};}
    \DataTypeTok{float}\NormalTok{ mass }\OperatorTok{=} \FloatTok{1.0}\BuiltInTok{f}\OperatorTok{;}
    \DataTypeTok{bool}\NormalTok{ fixed }\OperatorTok{=} \KeywordTok{false}\OperatorTok{;}
\OperatorTok{\};}
\end{Highlighting}
\end{Shaded}

\textbf{Units:} \texttt{position} m, \texttt{velocity} m/s,
\texttt{force} N, \texttt{mass} kg. A fixed node should not be
numerically integrated. The \texttt{force} field stores the \textbf{sum
of forces for the current timestep}, not a permanent applied force.

\subsection{\texorpdfstring{\texttt{Beam.hpp}: connected
spring/damper}{Beam.hpp: connected spring/damper}}\label{beam.hpp-connected-springdamper}

Recorded interface (field ordering may differ locally):

\begin{Shaded}
\begin{Highlighting}[]
\KeywordTok{struct}\NormalTok{ Beam }\OperatorTok{\{}
    \BuiltInTok{std::}\NormalTok{size\_t nodeA }\OperatorTok{=} \DecValTok{0}\OperatorTok{;}
    \BuiltInTok{std::}\NormalTok{size\_t nodeB }\OperatorTok{=} \DecValTok{0}\OperatorTok{;}
    \DataTypeTok{float}\NormalTok{ restLength }\OperatorTok{=} \FloatTok{0.0}\BuiltInTok{f}\OperatorTok{;}
    \DataTypeTok{float}\NormalTok{ stiffness }\OperatorTok{=} \FloatTok{1000.0}\BuiltInTok{f}\OperatorTok{;}
    \DataTypeTok{float}\NormalTok{ damping }\OperatorTok{=} \FloatTok{10.0}\BuiltInTok{f}\OperatorTok{;}
    \DataTypeTok{float}\NormalTok{ yieldStrain }\OperatorTok{=} \FloatTok{0.15}\BuiltInTok{f}\OperatorTok{;}
    \DataTypeTok{float}\NormalTok{ breakStrain }\OperatorTok{=} \FloatTok{0.60}\BuiltInTok{f}\OperatorTok{;}
    \DataTypeTok{float}\NormalTok{ plasticity }\OperatorTok{=} \FloatTok{2.0}\BuiltInTok{f}\OperatorTok{;}
    \DataTypeTok{bool}\NormalTok{ plasticDeformationEnabled }\OperatorTok{=} \KeywordTok{true}\OperatorTok{;}
    \DataTypeTok{bool}\NormalTok{ plasticallyDeformed }\OperatorTok{=} \KeywordTok{false}\OperatorTok{;}
    \DataTypeTok{bool}\NormalTok{ broken }\OperatorTok{=} \KeywordTok{false}\OperatorTok{;}
\OperatorTok{\};}
\end{Highlighting}
\end{Shaded}

The beam has no separate mass or bending DOF. \texttt{nodeA},
\texttt{nodeB} are \textbf{indices} into \texttt{world.nodes}.
\texttt{restLength} is the length at which the spring's elastic
component is zero. Later plastic deformation mutates that value, so
original-length-based comparisons would require a separate
original-length field.

\subsection{\texorpdfstring{\texttt{PhysicsWorld}: responsible for one
step}{PhysicsWorld: responsible for one step}}\label{physicsworld-responsible-for-one-step}

The previously recorded step sequence:

\begin{Shaded}
\begin{Highlighting}[]
\DataTypeTok{void}\NormalTok{ PhysicsWorld}\OperatorTok{::}\NormalTok{step}\OperatorTok{(}\DataTypeTok{float}\NormalTok{ dt}\OperatorTok{)} \OperatorTok{\{}
    \CommentTok{// timestep validation}
\NormalTok{    validateState}\OperatorTok{();}
\NormalTok{    resetForces}\OperatorTok{();}
\NormalTok{    applyExternalForces}\OperatorTok{();}
    \ControlFlowTok{if} \OperatorTok{(}\NormalTok{externalForceCallback}\OperatorTok{)}\NormalTok{ externalForceCallback}\OperatorTok{(*}\KeywordTok{this}\OperatorTok{);}
\NormalTok{    solveBeamForces}\OperatorTok{();}
\NormalTok{    integrate}\OperatorTok{(}\NormalTok{dt}\OperatorTok{);}
\NormalTok{    detectAndResolveGroundContacts}\OperatorTok{();}
\NormalTok{    updateDamage}\OperatorTok{(}\NormalTok{dt}\OperatorTok{);}
\NormalTok{    validateState}\OperatorTok{();}
\OperatorTok{\}}
\end{Highlighting}
\end{Shaded}

This is a \emph{recorded conceptual sequence}, not a byte-for-byte
snapshot of the newest local \texttt{PhysicsWorld.cpp}. Its order has
physical consequences: using forces from the current positions for
integration; applying ground corrections after integration; and
evaluating breakage after movement means a breaking beam may still have
contributed a force during the immediately preceding timestep.

\subsection{Why separate
responsibilities?}\label{why-separate-responsibilities}

\texttt{main.cpp} constructs a test \textbf{scenario}, processes UI
events and calls physics/render. \texttt{PhysicsWorld.cpp} evolves the
physical state. \texttt{Renderer.cpp} converts the latest state to GPU
draw commands. As a result, the same physics core could be called by
command-line tests, a 2D renderer, a 3D renderer or (later) a headless
benchmark.

\textbf{Mental model from Python:} \texttt{PhysicsWorld} is the domain
model, \texttt{Renderer} is a view, and \texttt{main.cpp} is the
controller/application entry point. Unlike a typical pure-Python
application, the C++ declarations, definitions, template instantiations
and linker all have to agree.

Read alongside source snapshots
\texttt{Source\ Snapshots/v0.4-GLAD/src/main.cpp} and
\texttt{.../Renderer.cpp}. See also 05 - Forces and Spring Dynamics and
15 - CMake Deep Dive.

\newpage

\section{Forces, Springs and Damping}\label{forces-springs-and-damping}

\subsection{Newton's second law}\label{newtons-second-law}

For a node with constant mass \(m\):

\[\sum \mathbf F = m\mathbf a, \qquad \mathbf a=\frac{\sum\mathbf F}{m}.\]

The world starts each timestep by clearing accumulated node forces, then
adding external and internal forces:

\begin{Shaded}
\begin{Highlighting}[]
\NormalTok{node}\OperatorTok{.}\NormalTok{force }\OperatorTok{=}\NormalTok{ Vector3D}\OperatorTok{\{}\FloatTok{0.0}\BuiltInTok{f}\OperatorTok{\};}
\ControlFlowTok{if} \OperatorTok{(!}\NormalTok{node}\OperatorTok{.}\NormalTok{fixed}\OperatorTok{)}\NormalTok{ node}\OperatorTok{.}\NormalTok{force }\OperatorTok{+=}\NormalTok{ node}\OperatorTok{.}\NormalTok{mass }\OperatorTok{*}\NormalTok{ gravity}\OperatorTok{;}
\end{Highlighting}
\end{Shaded}

If \(m=1\,\mathrm{kg}\) and \(g_y=-9.81\,\mathrm{m/s^2}\), the
gravitational force is \(F_y=-9.81\,\mathrm N\). A mass of 2 kg
experiences \(-19.62\,\mathrm N\), but \emph{the same gravitational
acceleration} when no other forces apply.

\subsection{The direction of a beam}\label{the-direction-of-a-beam}

For nodes A and B:

\[
\begin{aligned}
\mathbf d &= \mathbf x_B-\mathbf x_A,\\
L &= \lVert\mathbf d\rVert,\\
\hat{\mathbf d} &= \frac{\mathbf d}{L},\\
e &= L-L_0.
\end{aligned}
\]

Here \(L\) is current length, \(L_0\) is the \textbf{current rest
length}, and \(e\) is extension in metres. A very small \(L\) must be
guarded before dividing by it
(\texttt{if\ (length\ \textless{}\ 1e-6f)\ continue;}).

\subsection{Hooke's law plus axial
damping}\label{hookes-law-plus-axial-damping}

The force on node \textbf{A} in the project's implementation is

\[
\boxed{\mathbf F_A=\Big[k(L-L_0)+c\big((\mathbf v_B-\mathbf v_A)\cdot\hat{\mathbf d}\big)\Big]\hat{\mathbf d}}.
\]

The force on B is \(\mathbf F_B=-\mathbf F_A\) (Newton's third law).
Positive extension pulls A toward B. If B moves away from A along the
beam, the damping contribution opposes the relative motion by pulling A
toward B and B toward A. Verify the sign \emph{relative to the
endpoint}; changing the sign of damping is a common instability bug.

Recorded logic:

\begin{Shaded}
\begin{Highlighting}[]
\NormalTok{Vector3D delta }\OperatorTok{=}\NormalTok{ b}\OperatorTok{.}\NormalTok{position }\OperatorTok{{-}}\NormalTok{ a}\OperatorTok{.}\NormalTok{position}\OperatorTok{;}
\DataTypeTok{float}\NormalTok{ length }\OperatorTok{=}\NormalTok{ glm}\OperatorTok{::}\NormalTok{length}\OperatorTok{(}\NormalTok{delta}\OperatorTok{);}
\ControlFlowTok{if} \OperatorTok{(}\NormalTok{length }\OperatorTok{\textless{}} \FloatTok{1e{-}6}\BuiltInTok{f}\OperatorTok{)} \ControlFlowTok{continue}\OperatorTok{;}
\NormalTok{Vector3D direction }\OperatorTok{=}\NormalTok{ delta }\OperatorTok{/}\NormalTok{ length}\OperatorTok{;}
\DataTypeTok{float}\NormalTok{ extension }\OperatorTok{=}\NormalTok{ length }\OperatorTok{{-}}\NormalTok{ beam}\OperatorTok{.}\NormalTok{restLength}\OperatorTok{;}
\DataTypeTok{float}\NormalTok{ relativeVelocity }\OperatorTok{=}\NormalTok{ glm}\OperatorTok{::}\NormalTok{dot}\OperatorTok{(}\NormalTok{b}\OperatorTok{.}\NormalTok{velocity }\OperatorTok{{-}}\NormalTok{ a}\OperatorTok{.}\NormalTok{velocity}\OperatorTok{,}\NormalTok{ direction}\OperatorTok{);}
\DataTypeTok{float}\NormalTok{ springForce }\OperatorTok{=}\NormalTok{ beam}\OperatorTok{.}\NormalTok{stiffness }\OperatorTok{*}\NormalTok{ extension}\OperatorTok{;}
\DataTypeTok{float}\NormalTok{ dampingForce }\OperatorTok{=}\NormalTok{ beam}\OperatorTok{.}\NormalTok{damping }\OperatorTok{*}\NormalTok{ relativeVelocity}\OperatorTok{;}
\NormalTok{Vector3D force }\OperatorTok{=} \OperatorTok{(}\NormalTok{springForce }\OperatorTok{+}\NormalTok{ dampingForce}\OperatorTok{)} \OperatorTok{*}\NormalTok{ direction}\OperatorTok{;}
\ControlFlowTok{if} \OperatorTok{(!}\NormalTok{a}\OperatorTok{.}\NormalTok{fixed}\OperatorTok{)}\NormalTok{ a}\OperatorTok{.}\NormalTok{force }\OperatorTok{+=}\NormalTok{ force}\OperatorTok{;}
\ControlFlowTok{if} \OperatorTok{(!}\NormalTok{b}\OperatorTok{.}\NormalTok{fixed}\OperatorTok{)}\NormalTok{ b}\OperatorTok{.}\NormalTok{force }\OperatorTok{{-}=}\NormalTok{ force}\OperatorTok{;}
\end{Highlighting}
\end{Shaded}

\textbf{Units:} \(k\) N/m; \(c\) N·s/m; \(L-L_0\) m; \(\dot L\) m/s;
both resulting terms are N. These beams transmit \textbf{axial} load
only, not moments or direct bending. Cross-bracing creates network-level
resistance to shear.

\subsection{One-dimensional worked
calculation}\label{one-dimensional-worked-calculation}

Take \(L_0=1.00\) m, \(L=1.03\) m, \(k=1000\) N/m, \(c=10\) N·s/m, and
\(\dot L=0.2\) m/s:

\[F=k(0.03)+c(0.2)=30+2=32\,\mathrm N.\]

The magnitude is 32 N. Node A experiences \(+32\hat{\mathbf d}\) and
node B \(-32\hat{\mathbf d}\). If the beam is compressing, extension and
possibly axial velocity are negative, reversing the sign.

\subsection{Spring energy and damping}\label{spring-energy-and-damping}

An ideal undamped Hookean beam has potential energy:

\[U_s=\tfrac12 k(L-L_0)^2.\]

The axial damper dissipates energy at instantaneous rate

\[P_{d}=-c\dot L^2\le 0,\]

assuming the velocity projection and force signs above. This is a useful
test for damping implementation. \textbf{Caveat:} node ground clamping,
plastic changes to rest length and explicit integration alter the simple
conservative-energy budget.

\subsection{What is not modelled yet?}\label{what-is-not-modelled-yet}

\begin{itemize}
\tightlist
\item
  No rotation or moment constraints at beam endpoints.
\item
  No finite cross-section, distributed mass or bending stiffness for a
  single beam.
\item
  No continuous collision detection or self-collision.
\item
  No force-calibrated fatigue or crash material model.
\item
  No physically derived plastic yield surface; current thresholds are
  intentionally simplified.
\end{itemize}

See 09 - Plasticity, Fracture and Damage and 18 - Limits and Future
Roadmap. \href{https://www.beamng.com/game/about/physics/}{BeamNG's
node--beam explanation} is useful conceptual background but does not
imply Solver equivalence.

\newpage

\section{Numerical Integration, Stability and
Timesteps}\label{numerical-integration-stability-and-timesteps}

\subsection{Continuous versus discrete
motion}\label{continuous-versus-discrete-motion}

Newton's laws are differential equations:

\[\frac{d\mathbf x}{dt}=\mathbf v,\qquad \frac{d\mathbf v}{dt}=\mathbf a(\mathbf x,\mathbf v,t).\]

The CPU computes states at discrete times \(t_n\). SoloPhysics uses
\textbf{semi-implicit Euler}, also called symplectic Euler in this
context:

\[\boxed{\begin{aligned}\mathbf v_{n+1}&=\mathbf v_n+\Delta t\,\mathbf a_n,\\\mathbf x_{n+1}&=\mathbf x_n+\Delta t\,\mathbf v_{n+1}.\end{aligned}}\]

Contrast explicit Euler:
\(\mathbf x_{n+1}=\mathbf x_n+\Delta t\,\mathbf v_n\). Using the
\emph{new} velocity for the position step generally behaves better for
simple oscillators but does \textbf{not} make arbitrarily stiff springs
stable.

\begin{Shaded}
\begin{Highlighting}[]
\NormalTok{Vector3D acceleration }\OperatorTok{=}\NormalTok{ node}\OperatorTok{.}\NormalTok{force }\OperatorTok{/}\NormalTok{ node}\OperatorTok{.}\NormalTok{mass}\OperatorTok{;}
\NormalTok{node}\OperatorTok{.}\NormalTok{velocity }\OperatorTok{+=}\NormalTok{ acceleration }\OperatorTok{*}\NormalTok{ dt}\OperatorTok{;}
\NormalTok{node}\OperatorTok{.}\NormalTok{position }\OperatorTok{+=}\NormalTok{ node}\OperatorTok{.}\NormalTok{velocity }\OperatorTok{*}\NormalTok{ dt}\OperatorTok{;}
\end{Highlighting}
\end{Shaded}

A \texttt{fixed} node skips these updates. Nonpositive/NaN mass requires
validation before division.

\subsection{Choosing a small step}\label{choosing-a-small-step}

At \(\Delta t=1/240\) seconds, the simulation attempts 240 updates for
every simulated second. The single hanging-mass test initially used
\(1/120\) seconds; the later interactive application used \(1/240\). A
window rendering at 60 frames per second normally receives about four
physics steps per frame.

A rough scale for a 1 kg body on a 1000 N/m spring is

\[\omega_n=\sqrt{k/m}=\sqrt{1000}\approx31.62\,\mathrm{rad/s},\]

\[\omega_n\Delta t\approx0.132\quad(240\,\mathrm{Hz}).\]

Small \(\omega_n\Delta t\) generally improves numerical behaviour.
Coupled beams have multiple vibration modes, so one scalar spring
estimate is \textbf{not} a proof of overall stability. Damping too high
can also cause numerical trouble with explicit schemes.

\subsection{The frame-time
accumulator}\label{the-frame-time-accumulator}

Current v0.4 \texttt{main.cpp} includes:

\begin{Shaded}
\begin{Highlighting}[]
\KeywordTok{constexpr} \DataTypeTok{double}\NormalTok{ fixedDt }\OperatorTok{=} \FloatTok{1.0} \OperatorTok{/} \FloatTok{240.0}\OperatorTok{;}
\DataTypeTok{double}\NormalTok{ previousTime }\OperatorTok{=}\NormalTok{ glfwGetTime}\OperatorTok{();}
\DataTypeTok{double}\NormalTok{ accumulator }\OperatorTok{=} \FloatTok{0.0}\OperatorTok{;}

\ControlFlowTok{while} \OperatorTok{(!}\NormalTok{renderer}\OperatorTok{.}\NormalTok{shouldClose}\OperatorTok{())} \OperatorTok{\{}
    \DataTypeTok{double}\NormalTok{ now }\OperatorTok{=}\NormalTok{ glfwGetTime}\OperatorTok{();}
    \DataTypeTok{double}\NormalTok{ frameTime }\OperatorTok{=} \BuiltInTok{std::}\NormalTok{clamp}\OperatorTok{(}\NormalTok{now }\OperatorTok{{-}}\NormalTok{ previousTime}\OperatorTok{,} \FloatTok{0.0}\OperatorTok{,} \FloatTok{0.1}\OperatorTok{);}
\NormalTok{    previousTime }\OperatorTok{=}\NormalTok{ now}\OperatorTok{;}
\NormalTok{    accumulator }\OperatorTok{+=}\NormalTok{ frameTime}\OperatorTok{;}
    \ControlFlowTok{while} \OperatorTok{(}\NormalTok{accumulator }\OperatorTok{\textgreater{}=}\NormalTok{ fixedDt}\OperatorTok{)} \OperatorTok{\{}
\NormalTok{        world}\OperatorTok{.}\NormalTok{step}\OperatorTok{(}\KeywordTok{static\_cast}\OperatorTok{\textless{}}\DataTypeTok{float}\OperatorTok{\textgreater{}(}\NormalTok{fixedDt}\OperatorTok{));}
\NormalTok{        accumulator }\OperatorTok{{-}=}\NormalTok{ fixedDt}\OperatorTok{;}
    \OperatorTok{\}}
\NormalTok{    renderer}\OperatorTok{.}\NormalTok{render}\OperatorTok{(}\NormalTok{world}\OperatorTok{,}\NormalTok{ grabbedNode}\OperatorTok{,}\NormalTok{ hoveredNode}\OperatorTok{,}\NormalTok{ mouseTarget}\OperatorTok{);}
\OperatorTok{\}}
\end{Highlighting}
\end{Shaded}

This excerpt omits actual input processing, pause branch, and
diagnostics for clarity. The \emph{actual} source retains them.
\texttt{glfwGetTime()} returns elapsed time; \texttt{std::clamp} limits
a single frame to at most 0.1 s of physics catch-up. That means a
lengthy stall can cause simulated time to fall behind real elapsed time.
The accumulator preserves unsimulated fractional time within the chosen
frame-time cap.

\textbf{Important:} fixed-step does not guarantee fully deterministic
results across all hardware, floating-point modes or different
input-event timings. The app currently samples mouse input once per
render loop; it does not record/replay a timestamped input stream.

\subsection{Critical damping and
oscillations}\label{critical-damping-and-oscillations}

For an isolated 1D mass--spring--damper:

\[m\ddot x+c\dot x+kx=0,\]

\[\omega_n=\sqrt{\frac{k}{m}},\qquad c_{crit}=2\sqrt{km},\qquad\zeta=\frac{c}{2\sqrt{km}}.\]

\begin{itemize}
\tightlist
\item
  \(\zeta<1\): underdamped, oscillates while decaying.
\item
  \(\zeta=1\): critically damped in the \emph{ideal isolated system}.
\item
  \(\zeta>1\): overdamped.
\end{itemize}

For \(m=1\) kg, \(k=100\) N/m, \(c_{crit}=20\) N·s/m. With \(c=5\),
\(\zeta=0.25\); with \(c=0.6\), \(\zeta=0.03\). These idealised formulas
are excellent for validating the hanging-mass test but don't transfer
directly to every mode of a 28-beam network.

\subsection{Numerical checklist}\label{numerical-checklist}

\begin{enumerate}
\def\labelenumi{\arabic{enumi}.}
\tightlist
\item
  Validate finite and positive masses and finite positions/velocities.
\item
  Avoid normalising zero-length beams.
\item
  Monitor maximum speed/force/strain if nodes explode or vanish.
\item
  Apply forces \textbf{after} resetting accumulators.
\item
  Ensure each action uses a known fixed \(\Delta t\).
\item
  Use controlled tests with known answers before a complex vehicle.
\end{enumerate}

Additional reading: Glenn Fiedler,
\href{https://gafferongames.com/post/fix_your_timestep/}{Fix Your
Timestep!}. See 07 - The Hanging Mass Experiment.

\newpage

\section{First Verification: Suspended
Mass}\label{first-verification-suspended-mass}

\subsection{Experiment design}\label{experiment-design}

The first successful validation used a \textbf{single 1 kg mass}
suspended from a fixed anchor at \(y=3\) m by a beam of initial length
\(L_0=1\) m, stiffness \(k=100\) N/m and damping initially about \(5\)
N·s/m. Gravity pointed down with magnitude \(9.81\) m/s².

\begin{verbatim}
Anchor at y=3.0000 m   [fixed]
            |
            | beam, L0=1 m
            |
Mass initially y=2.0000 m [1 kg]
            ↓ gravity
\end{verbatim}

\subsection{\texorpdfstring{Predict the resting height \emph{before}
simulation}{Predict the resting height before simulation}}\label{predict-the-resting-height-before-simulation}

At equilibrium, the spring must exert \(mg=9.81\) N upward. The spring
therefore extends by

\[\Delta L=\frac{mg}{k}=\frac{1\times9.81}{100}=0.0981\,\mathrm m.\]

Its settled length is \(L=1.0981\) m. Because the mass is suspended
below the anchor,

\[\boxed{y_{eq}=3.0-1.0981=1.9019\,\mathrm m}.\]

The actual CLI output converged to approximately \textbf{1.9019 m},
matching this prediction at the reported precision. This gave initial
evidence that gravity, axial spring direction and integration were
coherently implemented.

\subsection{Damping experiment}\label{damping-experiment}

Keeping \(m=1\), \(k=100\) and changing damping from roughly \(c=5\) to
\(c=0.6\) caused oscillations to persist longer, as expected from
\(\zeta=c/(2\sqrt{km})\). A key experimental precaution was to disable
plastic deformation (\texttt{plasticDeformationEnabled=false}) and raise
\texttt{breakStrain} so the spring's rest length would not change during
this elastic baseline test.

\begin{quote}
\textbf{Note:} Why this is a strong engineering test A test is most
useful when you can calculate an expected result without relying on the
program. Otherwise, incorrect code may still draw a convincing
animation.
\end{quote}

\subsection{Lessons recorded}\label{lessons-recorded}

\begin{itemize}
\tightlist
\item
  Be sure to save \texttt{main.cpp} before rebuilding; an earlier test
  accidentally ran with unchanged source and misleadingly similar
  output.
\item
  Print damping, stiffness, plasticity state and rest length alongside
  results.
\item
  Compare analytical equilibrium, not just whether a mass visually stops
  moving.
\item
  Damping should affect oscillatory transient behaviour, not ideal
  static equilibrium while rest length is fixed.
\end{itemize}

\subsection{Suggested regression
assertion}\label{suggested-regression-assertion}

\begin{Shaded}
\begin{Highlighting}[]
\CommentTok{// Illustrative test only: assumes you set up the above world and}
\CommentTok{// run it long enough to settle without plasticity.}
\AttributeTok{const} \DataTypeTok{float}\NormalTok{ expectedY }\OperatorTok{=} \FloatTok{1.9019}\BuiltInTok{f}\OperatorTok{;}
\AttributeTok{const} \DataTypeTok{float}\NormalTok{ tolerance }\OperatorTok{=} \FloatTok{0.003}\BuiltInTok{f}\OperatorTok{;}
\OtherTok{assert}\OperatorTok{(}\BuiltInTok{std::}\NormalTok{abs}\OperatorTok{(}\NormalTok{world}\OperatorTok{.}\NormalTok{nodes}\OperatorTok{[}\NormalTok{massNodeIndex}\OperatorTok{].}\NormalTok{position}\OperatorTok{.}\NormalTok{y }\OperatorTok{{-}}\NormalTok{ expectedY}\OperatorTok{)} \OperatorTok{\textless{}}\NormalTok{ tolerance}\OperatorTok{);}
\end{Highlighting}
\end{Shaded}

This illustrative assertion is not an existing project test. It shows
how to automate the observation. See 17 - Verification Tests and
Exercises.

\newpage

\section{Square Drop, Impact and Resting-Contact
Fix}\label{square-drop-impact-and-resting-contact-fix}

\subsection{From one spring to a connected
structure}\label{from-one-spring-to-a-connected-structure}

The 2D test replaced the suspended node with four nodes placed in the XY
plane:

\begin{longtable}[]{@{}ll@{}}
\toprule\noalign{}
Node & Initial position (m) \\
\midrule\noalign{}
\endhead
\bottomrule\noalign{}
\endlastfoot
0, top left & \((-0.5,3,0)\) \\
1, top right & \((0.5,3,0)\) \\
2, bottom left & \((-0.5,2,0)\) \\
3, bottom right & \((0.5,2,0)\) \\
\end{longtable}

Six beams connected four perimeter edges and two crossing diagonals.
Diagonal cross-bracing stops the square becoming an unconstrained
pin-jointed parallelogram. With \(k=1000\) N/m, \(c=10\) N·s/m and each
node of mass 1 kg, the shape began falling from rest.

\begin{Shaded}
\begin{Highlighting}[]
\NormalTok{flowchart TB}
\NormalTok{ A((0)) {-}{-}{-} B((1))}
\NormalTok{ A {-}{-}{-} C((2))}
\NormalTok{ B {-}{-}{-} D((3))}
\NormalTok{ C {-}{-}{-} D}
\NormalTok{ A {-}{-}{-} D}
\NormalTok{ B {-}{-}{-} C}
\end{Highlighting}
\end{Shaded}

(The diagram is connectivity, not the actual screen positions.)

\subsection{Baseline free fall}\label{baseline-free-fall}

Before contact, all four nodes experienced almost identical downward
motion. At \(t=0.25\) s, the reported speed was approximately
\(v_y=-2.4525\) m/s, matching \(-gt=-9.81\times0.25\).

At \(t=0.50\) s, \(v_y=-4.9050\) m/s and bottom nodes were at
\(y=0.7635\) m. These are valuable baseline checks: internal beams
should not create a net external force on an initially undeformed freely
falling assembly.

The bottom starts at 2 m, so continuous ideal free-fall predicts first
contact near

\[t_{impact}=\sqrt{2h/g}=\sqrt{4/9.81}\approx0.639\,\mathrm s.\]

The logged data at 0.75 s already showed impact, rebound and
compression.

\subsection{The micro-bouncing bug}\label{the-micro-bouncing-bug}

The original ground solver clamped below-ground positions to \(y=0\) and
multiplied downward velocity by a restitution coefficient. At rest,
gravity repeatedly caused tiny penetration; the solver repeatedly
reversed velocity. Earlier results showed bottom nodes at \(y=0\) but
with persistent \(v_y\approx+0.0136\) m/s. That is an \textbf{artefact}
of the contact correction, not physical perpetual bouncing.

A simple resting threshold was introduced:

\begin{Shaded}
\begin{Highlighting}[]
\KeywordTok{constexpr} \DataTypeTok{float}\NormalTok{ restingVelocityThreshold }\OperatorTok{=} \FloatTok{0.2}\BuiltInTok{f}\OperatorTok{;}
\ControlFlowTok{if} \OperatorTok{(}\NormalTok{node}\OperatorTok{.}\NormalTok{position}\OperatorTok{.}\NormalTok{y }\OperatorTok{\textless{}}\NormalTok{ groundY}\OperatorTok{)} \OperatorTok{\{}
\NormalTok{    node}\OperatorTok{.}\NormalTok{position}\OperatorTok{.}\NormalTok{y }\OperatorTok{=}\NormalTok{ groundY}\OperatorTok{;}
    \ControlFlowTok{if} \OperatorTok{(}\NormalTok{node}\OperatorTok{.}\NormalTok{velocity}\OperatorTok{.}\NormalTok{y }\OperatorTok{\textless{}} \FloatTok{0.0}\BuiltInTok{f}\OperatorTok{)} \OperatorTok{\{}
        \DataTypeTok{float}\NormalTok{ impactSpeed }\OperatorTok{=} \OperatorTok{{-}}\NormalTok{node}\OperatorTok{.}\NormalTok{velocity}\OperatorTok{.}\NormalTok{y}\OperatorTok{;}
\NormalTok{        node}\OperatorTok{.}\NormalTok{velocity}\OperatorTok{.}\NormalTok{y }\OperatorTok{=}\NormalTok{ impactSpeed }\OperatorTok{\textless{}}\NormalTok{ restingVelocityThreshold}
                        \OperatorTok{?} \FloatTok{0.0}\BuiltInTok{f}
                        \OperatorTok{:} \OperatorTok{{-}}\NormalTok{node}\OperatorTok{.}\NormalTok{velocity}\OperatorTok{.}\NormalTok{y }\OperatorTok{*}\NormalTok{ restitution}\OperatorTok{;}
    \OperatorTok{\}}
\NormalTok{    node}\OperatorTok{.}\NormalTok{velocity}\OperatorTok{.}\NormalTok{x }\OperatorTok{*=} \OperatorTok{(}\FloatTok{1.0}\BuiltInTok{f} \OperatorTok{{-}}\NormalTok{ friction}\OperatorTok{);}
\NormalTok{    node}\OperatorTok{.}\NormalTok{velocity}\OperatorTok{.}\NormalTok{z }\OperatorTok{*=} \OperatorTok{(}\FloatTok{1.0}\BuiltInTok{f} \OperatorTok{{-}}\NormalTok{ friction}\OperatorTok{);}
\OperatorTok{\}}
\end{Highlighting}
\end{Shaded}

The displayed post-fix results settled to bottom \(y=0\), \(v_y=0\) (at
print precision) and upper \(y\approx0.9926\) m by approximately 2.5 s.
This was successful for the test.

\begin{quote}
\textbf{Note:} Not a physically complete contact solver This is a
position clamp and velocity heuristic. It does not compute persistent
normal impulses, static friction constraints, or continuous collision
detection. Multiplying tangential velocity by \texttt{(1\ -\ friction)}
each contact step is framerate/step-dependent and not Coulomb friction.
Do not treat it as a production crash solver.
\end{quote}

\subsection{Why the final upper height was below 1
m}\label{why-the-final-upper-height-was-below-1-m}

At rest, the bottom row sits at \(y=0\) while the upper row is around
\(0.9926\) m, slightly lower than the square's original 1 m height. The
network's axial springs are carrying gravitational loads and
compressing. The exact load distribution depends on beam topology and
constraints, so do not apply the single-hanging-mass formula blindly to
the square.

This experiment validated a multi-node topology, elastic loads, impact
response and the importance of separating numerical artefacts from
genuine material behaviour.

\newpage

\section{Plastic Deformation, Fracture and Damage
Flags}\label{plastic-deformation-fracture-and-damage-flags}

\subsection{The three regimes}\label{the-three-regimes}

\begin{enumerate}
\def\labelenumi{\arabic{enumi}.}
\tightlist
\item
  \textbf{Elastic:} beam stretches/compresses then returns toward its
  rest length when the load disappears.
\item
  \textbf{Plastic:} beam's rest length itself changes; deformation
  persists when forces stop.
\item
  \textbf{Fractured:} beam is marked broken and its force contribution
  is skipped thereafter.
\end{enumerate}

SoloPhysics implements a simplified \textbf{strain-threshold} model. For
current length \(L\) and current rest length \(L_0\),

\[\varepsilon=\frac{L-L_0}{L_0},\qquad \varepsilon_{abs}=|\varepsilon|.\]

When \(|\varepsilon|\ge\varepsilon_{break}\), \texttt{beam.broken=true}.
Otherwise, if plastic deformation is enabled and
\(|\varepsilon|>\varepsilon_y\), the code gradually moves \(L_0\) toward
a yield-boundary length.

\subsection{Recorded plastic update}\label{recorded-plastic-update}

\begin{Shaded}
\begin{Highlighting}[]
\DataTypeTok{float}\NormalTok{ factor }\OperatorTok{=} \BuiltInTok{std::}\NormalTok{clamp}\OperatorTok{(}\NormalTok{beam}\OperatorTok{.}\NormalTok{plasticity }\OperatorTok{*}\NormalTok{ dt}\OperatorTok{,} \FloatTok{0.0}\BuiltInTok{f}\OperatorTok{,} \FloatTok{1.0}\BuiltInTok{f}\OperatorTok{);}
\DataTypeTok{float}\NormalTok{ yieldBoundary }\OperatorTok{=}\NormalTok{ beam}\OperatorTok{.}\NormalTok{restLength }\OperatorTok{*}
    \OperatorTok{(}\FloatTok{1.0}\BuiltInTok{f} \OperatorTok{+} \BuiltInTok{std::}\NormalTok{copysign}\OperatorTok{(}\NormalTok{beam}\OperatorTok{.}\NormalTok{yieldStrain}\OperatorTok{,}\NormalTok{ strain}\OperatorTok{));}
\DataTypeTok{float}\NormalTok{ excess }\OperatorTok{=}\NormalTok{ length }\OperatorTok{{-}}\NormalTok{ yieldBoundary}\OperatorTok{;}
\DataTypeTok{float}\NormalTok{ oldRestLength }\OperatorTok{=}\NormalTok{ beam}\OperatorTok{.}\NormalTok{restLength}\OperatorTok{;}
\NormalTok{beam}\OperatorTok{.}\NormalTok{restLength }\OperatorTok{+=}\NormalTok{ excess }\OperatorTok{*}\NormalTok{ factor}\OperatorTok{;}
\NormalTok{beam}\OperatorTok{.}\NormalTok{restLength }\OperatorTok{=} \BuiltInTok{std::}\NormalTok{max}\OperatorTok{(}\NormalTok{beam}\OperatorTok{.}\NormalTok{restLength}\OperatorTok{,} \FloatTok{1e{-}6}\BuiltInTok{f}\OperatorTok{);}
\ControlFlowTok{if} \OperatorTok{(}\BuiltInTok{std::}\NormalTok{abs}\OperatorTok{(}\NormalTok{beam}\OperatorTok{.}\NormalTok{restLength }\OperatorTok{{-}}\NormalTok{ oldRestLength}\OperatorTok{)} \OperatorTok{\textgreater{}} \FloatTok{1e{-}7}\BuiltInTok{f}\OperatorTok{)}
\NormalTok{    beam}\OperatorTok{.}\NormalTok{plasticallyDeformed }\OperatorTok{=} \KeywordTok{true}\OperatorTok{;}
\end{Highlighting}
\end{Shaded}

This is a simplification recorded during the earlier 2D implementation;
verify the newest \texttt{PhysicsWorld.cpp} before changing it.
\texttt{std::copysign} keeps tension and compression yield boundaries on
the appropriate side of zero strain. \texttt{std::clamp} limits how far
the rest length can move during one step.

\subsection{Important limitation: moving strain
reference}\label{important-limitation-moving-strain-reference}

Because strain is divided by \textbf{the updated rest length}, yielding
can reduce current strain and can postpone fracture during slow pulling.
A more rigorous model should distinguish:

\[\varepsilon_{total}=\frac{L-L_{original}}{L_{original}},\qquad
\varepsilon_{elastic}=\frac{L-L_{plastic\,rest}}{L_{plastic\,rest}}\]

and track a history variable for accumulated plastic strain, damage or
dissipated plastic work. A stress/force-based yield model and material
calibration would be better for vehicle physics.

\subsection{Actual v0.4 demonstration
parameters}\label{actual-v0.4-demonstration-parameters}

From the checked 3D \texttt{src/main.cpp} bundle:

\begin{longtable}[]{@{}lrl@{}}
\toprule\noalign{}
Variable & Demo setting & Purpose \\
\midrule\noalign{}
\endhead
\bottomrule\noalign{}
\endlastfoot
\texttt{stiffness} & 650 N/m & Elastic axial response \\
\texttt{damping} & 12 N·s/m & Axial viscous damping \\
\texttt{yieldStrain} & 0.06 & Plastic update triggers above 6\% \\
\texttt{breakStrain} & 0.30 & Beam fractures above 30\% current-rest
strain \\
\texttt{plasticity} & 0.6 & Rest-length adaptation rate \\
\texttt{mouseStiffness} & 500 N/m & User's virtual drag spring \\
\texttt{mouseDamping} & 35 N·s/m & Drag force damper \\
\texttt{maxMouseForce} & 1500 N & Clamp for large mouse pulls \\
\end{longtable}

These \textbf{aren't validated properties of steel, aluminium or vehicle
parts}; they were chosen for an interactive demo.

\subsection{Colours and state
priority}\label{colours-and-state-priority}

The GLAD renderer uses green for near rest length, blue for compression,
red for extension, orange for a plastically deformed beam, and grey for
broken. Plastic history takes priority over current elastic strain (so
an orange beam can also currently be stretched or compressed). Broken
beams are still drawn as grey straight lines between endpoints
\textbf{for debugging}, not as dangling fragments.

\subsection{First damage debugging
incident}\label{first-damage-debugging-incident}

At one point, the square visibly responded to dragging but never yielded
or broke because \texttt{main.cpp} still contained:

\begin{Shaded}
\begin{Highlighting}[]
\NormalTok{beam}\OperatorTok{.}\NormalTok{plasticDeformationEnabled }\OperatorTok{=} \KeywordTok{false}\OperatorTok{;}
\NormalTok{beam}\OperatorTok{.}\NormalTok{breakStrain }\OperatorTok{=} \FloatTok{100.0}\BuiltInTok{f}\OperatorTok{;}
\end{Highlighting}
\end{Shaded}

The fix was to enable yielding, set feasible demonstration thresholds,
anchor one node to force actual structural strain, and print
\texttt{Plastic\ beams}/\texttt{Broken\ beams} alongside maximum strain.
The user later confirmed the complete interactive damage test worked.

See 16 - Errors, Diagnostics and Lessons and 18 - Limits and Future
Roadmap.

\newpage

\section{The First 2D GLFW/OpenGL
Renderer}\label{the-first-2d-glfwopengl-renderer}

\subsection{Why add graphics only after headless
tests?}\label{why-add-graphics-only-after-headless-tests}

The CLI square-drop already produced useful positions and velocities. A
renderer made the dynamics understandable visually, but it shouldn't
change the physics results. The first GLFW version received
\texttt{const\ PhysicsWorld\&} and used an \textbf{orthographic XY
projection} of a 3D physics state.

\subsection{Basic GLFW lifecycle}\label{basic-glfw-lifecycle}

\begin{Shaded}
\begin{Highlighting}[]
\ControlFlowTok{if} \OperatorTok{(!}\NormalTok{glfwInit}\OperatorTok{())} \ControlFlowTok{throw} \BuiltInTok{std::}\NormalTok{runtime\_error}\OperatorTok{(}\StringTok{"GLFW init failed"}\OperatorTok{);}
\NormalTok{GLFWwindow}\OperatorTok{*}\NormalTok{ window }\OperatorTok{=}\NormalTok{ glfwCreateWindow}\OperatorTok{(}\DecValTok{1000}\OperatorTok{,} \DecValTok{700}\OperatorTok{,} \StringTok{"SoloPhysics"}\OperatorTok{,} \KeywordTok{nullptr}\OperatorTok{,} \KeywordTok{nullptr}\OperatorTok{);}
\ControlFlowTok{if} \OperatorTok{(!}\NormalTok{window}\OperatorTok{)} \ControlFlowTok{throw} \BuiltInTok{std::}\NormalTok{runtime\_error}\OperatorTok{(}\StringTok{"Window creation failed"}\OperatorTok{);}
\NormalTok{glfwMakeContextCurrent}\OperatorTok{(}\NormalTok{window}\OperatorTok{);}
\NormalTok{glfwSwapInterval}\OperatorTok{(}\DecValTok{1}\OperatorTok{);}
\ControlFlowTok{while} \OperatorTok{(!}\NormalTok{glfwWindowShouldClose}\OperatorTok{(}\NormalTok{window}\OperatorTok{))} \OperatorTok{\{}
\NormalTok{    glfwPollEvents}\OperatorTok{();}
    \CommentTok{// step physics and draw}
\NormalTok{    glfwSwapBuffers}\OperatorTok{(}\NormalTok{window}\OperatorTok{);}
\OperatorTok{\}}
\NormalTok{glfwDestroyWindow}\OperatorTok{(}\NormalTok{window}\OperatorTok{);}
\NormalTok{glfwTerminate}\OperatorTok{();}
\end{Highlighting}
\end{Shaded}

This is an illustrative lifecycle excerpt, \textbf{not} the complete
current code. GLFW creates native windows and OpenGL contexts; it does
not itself implement the physics engine or GPU shaders.
\texttt{glfwSwapInterval(1)} requests vsync; actual frame rate depends
on environment.

\subsection{Legacy OpenGL prototype}\label{legacy-opengl-prototype}

The original \texttt{Renderer.cpp} used a requested \textbf{OpenGL 2.1
compatibility context}, \texttt{glBegin(GL\_LINES)},
\texttt{glVertex2f(...)}, \texttt{glEnd()}, \texttt{glMatrixMode}, and
\texttt{glOrtho}. This was an intentionally short teaching step, not the
final architecture. Each frame drew:

\begin{itemize}
\tightlist
\item
  ground line;
\item
  all unbroken beams as coloured segments;
\item
  nodes as points.
\end{itemize}

\begin{Shaded}
\begin{Highlighting}[]
\NormalTok{glBegin}\OperatorTok{(}\NormalTok{GL\_LINES}\OperatorTok{);}
\ControlFlowTok{for} \OperatorTok{(}\AttributeTok{const}\NormalTok{ Beam}\OperatorTok{\&}\NormalTok{ beam }\OperatorTok{:}\NormalTok{ world}\OperatorTok{.}\NormalTok{beams}\OperatorTok{)} \OperatorTok{\{}
    \ControlFlowTok{if} \OperatorTok{(}\NormalTok{beam}\OperatorTok{.}\NormalTok{broken}\OperatorTok{)} \ControlFlowTok{continue}\OperatorTok{;}
    \AttributeTok{const}\NormalTok{ Node}\OperatorTok{\&}\NormalTok{ a }\OperatorTok{=}\NormalTok{ world}\OperatorTok{.}\NormalTok{nodes}\OperatorTok{[}\NormalTok{beam}\OperatorTok{.}\NormalTok{nodeA}\OperatorTok{];}
    \AttributeTok{const}\NormalTok{ Node}\OperatorTok{\&}\NormalTok{ b }\OperatorTok{=}\NormalTok{ world}\OperatorTok{.}\NormalTok{nodes}\OperatorTok{[}\NormalTok{beam}\OperatorTok{.}\NormalTok{nodeB}\OperatorTok{];}
\NormalTok{    glVertex2f}\OperatorTok{(}\NormalTok{a}\OperatorTok{.}\NormalTok{position}\OperatorTok{.}\NormalTok{x}\OperatorTok{,}\NormalTok{ a}\OperatorTok{.}\NormalTok{position}\OperatorTok{.}\NormalTok{y}\OperatorTok{);}
\NormalTok{    glVertex2f}\OperatorTok{(}\NormalTok{b}\OperatorTok{.}\NormalTok{position}\OperatorTok{.}\NormalTok{x}\OperatorTok{,}\NormalTok{ b}\OperatorTok{.}\NormalTok{position}\OperatorTok{.}\NormalTok{y}\OperatorTok{);}
\OperatorTok{\}}
\NormalTok{glEnd}\OperatorTok{();}
\end{Highlighting}
\end{Shaded}

\textbf{Syntax:} \texttt{glBegin} and \texttt{glEnd} are function calls.
\texttt{GL\_LINES} is a named API constant. The
\texttt{for\ (const\ Beam\&\ beam\ :\ ...)} loop accesses each beam by
read-only reference. The renderer reads node positions and translates
them into screen geometry.

\subsection{The early screen-space/world-space
trap}\label{the-early-screen-spaceworld-space-trap}

World position \(x,y\) is measured in metres; mouse location is measured
in window pixels. A screen-centred orthographic camera with vertical
span \(H\) converts cursor coordinates approximately as:

\[x_w=x_c+\left(2\frac{x_p}{W}-1\right)\frac{H}{2}\frac{W}{H_{pixel}},\]

\[y_w=y_c+\left(1-2\frac{y_p}{H_{pixel}}\right)\frac{H}{2}.\]

The actual prototype also handled framebuffer size separately from
window size for scaling on HiDPI setups. Its zoom and panning affected
the mapping.

\subsection{Why replace the renderer?}\label{why-replace-the-renderer}

Modern OpenGL core profiles intentionally omit legacy immediate-mode
drawing calls. Moving to OpenGL 3.3 shaders and vertex buffers makes the
engine more portable to modern contexts and introduces the GPU pipeline
required for larger meshes and later vehicle rendering.

See 13 - Modern OpenGL and GLSL and
\href{https://www.glfw.org/docs/latest/quick_guide.html}{GLFW's official
getting-started guide}.

\newpage

\section{Interactive Controls, Picking and Mouse
Springs}\label{interactive-controls-picking-and-mouse-springs}

\subsection{The problem: force accumulators were
reset}\label{the-problem-force-accumulators-were-reset}

The first drag attempt failed because the application added forces at
the wrong point in the update sequence. Each
\texttt{PhysicsWorld::step()} begins by setting all \texttt{node.force}
fields to zero. If \texttt{main.cpp} writes \texttt{node.force\ +=\ ...}
\textbf{before} \texttt{step()}, that force disappears on the very next
reset.

\textbf{Solution:} \texttt{PhysicsWorld} exposes an
\texttt{externalForceCallback} invoked \textbf{after} resetting forces
and applying gravity but \textbf{before} solving beams and integrating.

\begin{Shaded}
\begin{Highlighting}[]
\CommentTok{// Recorded interface}
\BuiltInTok{std::}\NormalTok{function}\OperatorTok{\textless{}}\DataTypeTok{void}\OperatorTok{(}\NormalTok{PhysicsWorld}\OperatorTok{\&)\textgreater{}}\NormalTok{ externalForceCallback}\OperatorTok{;}

\CommentTok{// In step() (illustrative excerpt)}
\NormalTok{resetForces}\OperatorTok{();}
\NormalTok{applyExternalForces}\OperatorTok{();}
\ControlFlowTok{if} \OperatorTok{(}\NormalTok{externalForceCallback}\OperatorTok{)}\NormalTok{ externalForceCallback}\OperatorTok{(*}\KeywordTok{this}\OperatorTok{);}
\NormalTok{solveBeamForces}\OperatorTok{();}
\NormalTok{integrate}\OperatorTok{(}\NormalTok{dt}\OperatorTok{);}
\end{Highlighting}
\end{Shaded}

\texttt{std::function\textless{}void(PhysicsWorld\&)\textgreater{}}
means: an object callable with one mutable \texttt{PhysicsWorld\&} and
returning \texttt{void}. Python analogy: a stored callable attribute
such as \texttt{world.external\_force\_callback\ =\ callback}.

\subsection{The virtual mouse spring}\label{the-virtual-mouse-spring}

\[\boxed{\mathbf F_{mouse}=k_m(\mathbf x_{target}-\mathbf x_{node})-c_m\mathbf v_{node}}.\]

The current 3D source contains:

\begin{Shaded}
\begin{Highlighting}[]
\KeywordTok{auto}\NormalTok{ applyMouseForce }\OperatorTok{=} \OperatorTok{[\&](}\NormalTok{PhysicsWorld}\OperatorTok{\&}\NormalTok{ physics}\OperatorTok{)} \OperatorTok{\{}
    \ControlFlowTok{if} \OperatorTok{(!}\NormalTok{grabbedNode }\OperatorTok{||} \OperatorTok{*}\NormalTok{grabbedNode }\OperatorTok{\textgreater{}=}\NormalTok{ physics}\OperatorTok{.}\NormalTok{nodes}\OperatorTok{.}\NormalTok{size}\OperatorTok{())} \ControlFlowTok{return}\OperatorTok{;}
\NormalTok{    Node}\OperatorTok{\&}\NormalTok{ node }\OperatorTok{=}\NormalTok{ physics}\OperatorTok{.}\NormalTok{nodes}\OperatorTok{[*}\NormalTok{grabbedNode}\OperatorTok{];}
    \ControlFlowTok{if} \OperatorTok{(}\NormalTok{node}\OperatorTok{.}\NormalTok{fixed}\OperatorTok{)} \ControlFlowTok{return}\OperatorTok{;}
\NormalTok{    Vector3D force }\OperatorTok{=}\NormalTok{ mouseStiffness }\OperatorTok{*} \OperatorTok{(}\NormalTok{mouseTarget }\OperatorTok{{-}}\NormalTok{ node}\OperatorTok{.}\NormalTok{position}\OperatorTok{)}
                   \OperatorTok{{-}}\NormalTok{ mouseDamping }\OperatorTok{*}\NormalTok{ node}\OperatorTok{.}\NormalTok{velocity}\OperatorTok{;}
    \KeywordTok{constexpr} \DataTypeTok{float}\NormalTok{ maxMouseForce }\OperatorTok{=} \FloatTok{1500.0}\BuiltInTok{f}\OperatorTok{;}
    \DataTypeTok{float}\NormalTok{ magnitude }\OperatorTok{=}\NormalTok{ glm}\OperatorTok{::}\NormalTok{length}\OperatorTok{(}\NormalTok{force}\OperatorTok{);}
    \ControlFlowTok{if} \OperatorTok{(}\NormalTok{magnitude }\OperatorTok{\textgreater{}}\NormalTok{ maxMouseForce}\OperatorTok{)}
\NormalTok{        force }\OperatorTok{*=}\NormalTok{ maxMouseForce }\OperatorTok{/}\NormalTok{ magnitude}\OperatorTok{;}
\NormalTok{    node}\OperatorTok{.}\NormalTok{force }\OperatorTok{+=}\NormalTok{ force}\OperatorTok{;}
\OperatorTok{\};}
\end{Highlighting}
\end{Shaded}

This code \textbf{does not teleport a node}; it applies a bounded spring
force, allowing the beam network to respond. \texttt{node.force\ +=}
adds to rather than replaces gravity and other external forces.

\subsection{Picking: from metres to
pixels}\label{picking-from-metres-to-pixels}

The original 2D picking selected the nearest node inside a fixed
world-space radius (initially 0.2 m, later enlarged for diagnosis). This
felt inconsistent as zoom changed.

The current 3D renderer uses
\texttt{glm::project(world.nodes{[}i{]}.position,\ viewMatrix(),\ projectionMatrix(),\ viewport)}
to transform each node into framebuffer coordinates. It compares
distance to the mouse cursor \textbf{in pixels} (19 px in the supplied
v0.4 caller), so the selection radius is consistent with screen scale.

\begin{Shaded}
\begin{Highlighting}[]
\KeywordTok{auto}\NormalTok{ hoveredNode }\OperatorTok{=}\NormalTok{ renderer}\OperatorTok{.}\NormalTok{pickNode}\OperatorTok{(}\NormalTok{world}\OperatorTok{,} \FloatTok{19.0}\BuiltInTok{f}\OperatorTok{);}
\ControlFlowTok{if} \OperatorTok{(}\NormalTok{left }\OperatorTok{\&\&} \OperatorTok{!}\NormalTok{previousLeft}\OperatorTok{)} \OperatorTok{\{}
\NormalTok{    grabbedNode }\OperatorTok{=}\NormalTok{ hoveredNode}\OperatorTok{;}
    \ControlFlowTok{if} \OperatorTok{(}\NormalTok{grabbedNode }\OperatorTok{\&\&}\NormalTok{ world}\OperatorTok{.}\NormalTok{nodes}\OperatorTok{[*}\NormalTok{grabbedNode}\OperatorTok{].}\NormalTok{fixed}\OperatorTok{)}
\NormalTok{        grabbedNode}\OperatorTok{.}\NormalTok{reset}\OperatorTok{();}
    \ControlFlowTok{if} \OperatorTok{(}\NormalTok{grabbedNode}\OperatorTok{)}
\NormalTok{        renderer}\OperatorTok{.}\NormalTok{beginDrag}\OperatorTok{(}\NormalTok{world}\OperatorTok{.}\NormalTok{nodes}\OperatorTok{[*}\NormalTok{grabbedNode}\OperatorTok{].}\NormalTok{position}\OperatorTok{);}
\OperatorTok{\}}
\end{Highlighting}
\end{Shaded}

\texttt{left\ \&\&\ !previousLeft} is an \textbf{edge detector}: act
once when the button transitions from released to pressed. Without it,
selection could restart every frame while held.

\subsection{3D dragging is not a 2D screen
coordinate}\label{d-dragging-is-not-a-2d-screen-coordinate}

In 3D, the cursor represents a \textbf{ray} extending out from the
camera. The program unprojects the near/far cursor points, forms a ray,
then intersects it with a plane passing through the selected node and
facing the camera:

\[\mathbf r(t)=\mathbf o+t\mathbf d, \quad \mathbf n\cdot(\mathbf r(t)-\mathbf p_0)=0,\]

\[\boxed{t=\frac{\mathbf n\cdot(\mathbf p_0-\mathbf o)}{\mathbf n\cdot\mathbf d}}.\]

If the denominator is close to zero, the ray is parallel to the plane;
return no target (\texttt{std::nullopt}). This gives intuitive movement
approximately parallel to the screen at the node's initial depth. It
does not implement arbitrary depth-axis gizmos yet.

\subsection{UI feedback}\label{ui-feedback}

\begin{itemize}
\tightlist
\item
  Normal node: yellow; fixed node: purple.
\item
  Hovered: cyan; grabbed: orange.
\item
  A yellowish drag-spring line joins grabbed node and cursor target.
\item
  \texttt{Space}, \texttt{R}, \texttt{G} and later \texttt{F} toggle
  pause, reset, gravity and the anchor; \texttt{Esc} exits.
\end{itemize}

\textbf{Debugging lesson:} seeing \texttt{Left\ mouse\ clicked}
confirmed GLFW, but wasn't proof that node picking succeeded; only
seeing selected IDs and applying forces confirmed the chain.

Further reading: \href{https://www.glfw.org/docs/latest/input.html}{GLFW
Input Guide}. See 03 - C++ for a Python Developer for lambdas and
optional values.

\newpage

\section{From Square to 3D Cube and Orbit
Camera}\label{from-square-to-3d-cube-and-orbit-camera}

\subsection{The v0.4 geometry}\label{the-v0.4-geometry}

The actual supplied v0.4 \texttt{main.cpp} creates \textbf{8 nodes},
four at \(y=2\) m and four at \(y=3\) m, with \(x,z\in\{-0.5,+0.5\}\) m.
Node 0 starts fixed. Gravity starts off, so the user can pull the
structure without free fall; pressing \texttt{F} releases the anchor,
and \texttt{G} enables gravity.

The edges form a fully connected 8-node network, \(K_8\), with

\[\binom{8}{2}=\frac{8\times7}{2}=28\text{ unique connections}.\]

The source enumerates 12 cube edges, 12 face diagonals and 4 body
diagonals. This is deliberately heavily braced; it is not a realistic
vehicle-body finite element model.

\begin{Shaded}
\begin{Highlighting}[]
\AttributeTok{const} \BuiltInTok{std::}\NormalTok{array}\OperatorTok{\textless{}}\NormalTok{Vector3D}\OperatorTok{,} \DecValTok{8}\OperatorTok{\textgreater{}}\NormalTok{ positions}\OperatorTok{\{\{}
    \OperatorTok{\{{-}}\FloatTok{0.5}\BuiltInTok{f}\OperatorTok{,} \FloatTok{2.0}\BuiltInTok{f}\OperatorTok{,} \OperatorTok{{-}}\FloatTok{0.5}\BuiltInTok{f}\OperatorTok{\},}
    \OperatorTok{\{} \FloatTok{0.5}\BuiltInTok{f}\OperatorTok{,} \FloatTok{2.0}\BuiltInTok{f}\OperatorTok{,} \OperatorTok{{-}}\FloatTok{0.5}\BuiltInTok{f}\OperatorTok{\},}
    \OperatorTok{\{} \FloatTok{0.5}\BuiltInTok{f}\OperatorTok{,} \FloatTok{2.0}\BuiltInTok{f}\OperatorTok{,}  \FloatTok{0.5}\BuiltInTok{f}\OperatorTok{\},}
    \OperatorTok{\{{-}}\FloatTok{0.5}\BuiltInTok{f}\OperatorTok{,} \FloatTok{2.0}\BuiltInTok{f}\OperatorTok{,}  \FloatTok{0.5}\BuiltInTok{f}\OperatorTok{\},}
    \OperatorTok{\{{-}}\FloatTok{0.5}\BuiltInTok{f}\OperatorTok{,} \FloatTok{3.0}\BuiltInTok{f}\OperatorTok{,} \OperatorTok{{-}}\FloatTok{0.5}\BuiltInTok{f}\OperatorTok{\},}
    \OperatorTok{\{} \FloatTok{0.5}\BuiltInTok{f}\OperatorTok{,} \FloatTok{3.0}\BuiltInTok{f}\OperatorTok{,} \OperatorTok{{-}}\FloatTok{0.5}\BuiltInTok{f}\OperatorTok{\},}
    \OperatorTok{\{} \FloatTok{0.5}\BuiltInTok{f}\OperatorTok{,} \FloatTok{3.0}\BuiltInTok{f}\OperatorTok{,}  \FloatTok{0.5}\BuiltInTok{f}\OperatorTok{\},}
    \OperatorTok{\{{-}}\FloatTok{0.5}\BuiltInTok{f}\OperatorTok{,} \FloatTok{3.0}\BuiltInTok{f}\OperatorTok{,}  \FloatTok{0.5}\BuiltInTok{f}\OperatorTok{\}}
\OperatorTok{\}\};}
\end{Highlighting}
\end{Shaded}

\texttt{std::array\textless{}Vector3D,8\textgreater{}} has exactly eight
GLM vectors. The \emph{two sets} of braces are common in
\texttt{std::array} aggregate initialisation and should not be confused
with a Python nested tuple of arbitrary size.

\subsection{Orbit camera geometry}\label{orbit-camera-geometry}

The camera stores a \textbf{target} (point to look at), a \textbf{yaw}
(left/right rotation), a \textbf{pitch} (up/down rotation), and a
\textbf{distance}. Given camera target \(\mathbf t\), distance \(r\),
yaw \(\theta\) and pitch \(\phi\):

\[\mathbf x_{camera}=\mathbf t+r\begin{bmatrix}
\cos\phi\sin\theta\\
\sin\phi\\
\cos\phi\cos\theta
\end{bmatrix}.\]

The renderer calls
\texttt{glm::lookAt(cameraPosition(),\ cameraTarget,\ \{0,1,0\})} to
obtain a view matrix. Right-drag changes yaw/pitch, middle-drag
translates the target, scrolling scales distance (clamped between 1.5
and 40 in the supplied version).

\subsection{From world coordinates to GPU
pixels}\label{from-world-coordinates-to-gpu-pixels}

A 3D point in homogeneous coordinates \(\mathbf p=(x,y,z,1)\) passes
through

\[\boxed{\mathbf p_{clip}=PVM\mathbf p_{model}}.\]

Current app geometry is directly in world coordinates, so its model
transform \(M\) is effectively identity and the code computes
\texttt{mvp\ =\ projectionMatrix()\ *\ viewMatrix()}. The matrix product
order matters: for column-vector convention, view acts first and
projection acts second.

The current projection is a perspective camera:

\begin{Shaded}
\begin{Highlighting}[]
\NormalTok{glm}\OperatorTok{::}\NormalTok{perspective}\OperatorTok{(}\NormalTok{glm}\OperatorTok{::}\NormalTok{radians}\OperatorTok{(}\FloatTok{55.0}\BuiltInTok{f}\OperatorTok{),}\NormalTok{ aspect}\OperatorTok{,} \FloatTok{0.05}\BuiltInTok{f}\OperatorTok{,} \FloatTok{100.0}\BuiltInTok{f}\OperatorTok{);}
\end{Highlighting}
\end{Shaded}

\texttt{55°} is vertical field of view;
\texttt{aspect\ =\ framebufferWidth/framebufferHeight}; \texttt{0.05} m
and \texttt{100} m set near/far clipping planes. Keeping the near plane
above zero avoids an invalid perspective projection.

\subsection{Screen-space picking and
unprojection}\label{screen-space-picking-and-unprojection}

\begin{itemize}
\tightlist
\item
  \texttt{glm::project}: world point → framebuffer X/Y/depth.
\item
  \texttt{glm::unProject}: framebuffer X/Y/depth → world point.
\item
  \texttt{cursorRay()}: two unprojected endpoints → a ray.
\item
  \texttt{beginDrag()}: record a point and a camera-facing plane normal.
\item
  \texttt{mouseDragTarget()}: ray/plane intersection.
\end{itemize}

This is \textbf{selection by projected point radius}, not mesh triangle
ray-intersection. It is appropriate for a node-based prototype but would
need occlusion, depth sorting and hit tests for complex visible
geometry.

\subsection{Scope of 3D collisions}\label{scope-of-3d-collisions}

The existing solver checks node height versus a flat \(y=groundY\)
plane. The renderer draws a \textbf{grid on the XZ plane}; the grid does
not itself simulate collision. The simple ground solver handles point
contacts, not edge/triangle collisions, contact patches, angular
friction or full cube--cube intersection.

See 13 - Modern OpenGL and GLSL and 18 - Limits and Future Roadmap.

\newpage

\section{Modern OpenGL, GPU Resources and
GLSL}\label{modern-opengl-gpu-resources-and-glsl}

\subsection{Why the OpenGL 3.3 change
matters}\label{why-the-opengl-3.3-change-matters}

OpenGL 2.1 compatibility mode allowed simple immediate-mode calls such
as \texttt{glBegin}. OpenGL 3.3 \textbf{core} instead requires defining
shaders and vertex input explicitly. The GPU draws the data you upload
into buffers. The current v0.4 renderer in the supplied bundle uses
\textbf{GLAD 2} to obtain function pointers and GLFW to
create/context-manage the window.

\subsection{Two shader stages}\label{two-shader-stages}

Actual vertex-shader logic in the supplied source:

\begin{Shaded}
\begin{Highlighting}[]
\PreprocessorTok{\#version 330 core}
\KeywordTok{layout}\OperatorTok{(}\DataTypeTok{location} \OperatorTok{=} \DecValTok{0}\OperatorTok{)} \DataTypeTok{in} \DataTypeTok{vec3}\NormalTok{ aPosition}\OperatorTok{;}
\KeywordTok{layout}\OperatorTok{(}\DataTypeTok{location} \OperatorTok{=} \DecValTok{1}\OperatorTok{)} \DataTypeTok{in} \DataTypeTok{vec3}\NormalTok{ aColour}\OperatorTok{;}
\KeywordTok{uniform} \DataTypeTok{mat4}\NormalTok{ uMVP}\OperatorTok{;}
\DataTypeTok{out} \DataTypeTok{vec3}\NormalTok{ vColour}\OperatorTok{;}
\DataTypeTok{void} \FunctionTok{main}\OperatorTok{()} \OperatorTok{\{}
    \BuiltInTok{gl\_Position} \OperatorTok{=}\NormalTok{ uMVP }\OperatorTok{*} \DataTypeTok{vec4}\OperatorTok{(}\NormalTok{aPosition}\OperatorTok{,} \FloatTok{1.0}\OperatorTok{);}
\NormalTok{    vColour }\OperatorTok{=}\NormalTok{ aColour}\OperatorTok{;}
\OperatorTok{\}}
\end{Highlighting}
\end{Shaded}

\texttt{layout(location\ =\ 0)} binds attribute index 0 to a 3-component
vertex position. \texttt{uniform} means one matrix value used across the
draw call; \texttt{out} passes interpolated data to the next shader
stage. \texttt{vec4(aPosition,1.0)} adds homogeneous coordinate \(w=1\)
so translations can be represented by a matrix multiplication.

Actual fragment-shader logic:

\begin{Shaded}
\begin{Highlighting}[]
\PreprocessorTok{\#version 330 core}
\DataTypeTok{in} \DataTypeTok{vec3}\NormalTok{ vColour}\OperatorTok{;}
\DataTypeTok{out} \DataTypeTok{vec4}\NormalTok{ FragColour}\OperatorTok{;}
\KeywordTok{uniform} \DataTypeTok{bool}\NormalTok{ uRoundPoints}\OperatorTok{;}
\DataTypeTok{void} \FunctionTok{main}\OperatorTok{()} \OperatorTok{\{}
    \KeywordTok{if} \OperatorTok{(}\NormalTok{uRoundPoints}\OperatorTok{)} \OperatorTok{\{}
        \DataTypeTok{vec2}\NormalTok{ p }\OperatorTok{=} \BuiltInTok{gl\_PointCoord} \OperatorTok{*} \FloatTok{2.0} \OperatorTok{{-}} \FloatTok{1.0}\OperatorTok{;}
        \KeywordTok{if} \OperatorTok{(}\BuiltInTok{dot}\OperatorTok{(}\NormalTok{p}\OperatorTok{,}\NormalTok{ p}\OperatorTok{)} \OperatorTok{\textgreater{}} \FloatTok{1.0}\OperatorTok{)} \KeywordTok{discard}\OperatorTok{;}
    \OperatorTok{\}}
\NormalTok{    FragColour }\OperatorTok{=} \DataTypeTok{vec4}\OperatorTok{(}\NormalTok{vColour}\OperatorTok{,} \FloatTok{1.0}\OperatorTok{);}
\OperatorTok{\}}
\end{Highlighting}
\end{Shaded}

For points, \texttt{gl\_PointCoord} gives a coordinate across the point
primitive. Pixels outside a unit circle are discarded, creating round
point sprites. \texttt{vec4(...,1.0)} assigns fully opaque alpha.

\subsection{Compile and link shaders}\label{compile-and-link-shaders}

\begin{Shaded}
\begin{Highlighting}[]
\NormalTok{GLuint shader }\OperatorTok{=}\NormalTok{ glCreateShader}\OperatorTok{(}\NormalTok{type}\OperatorTok{);}
\NormalTok{glShaderSource}\OperatorTok{(}\NormalTok{shader}\OperatorTok{,} \DecValTok{1}\OperatorTok{,} \OperatorTok{\&}\NormalTok{source}\OperatorTok{,} \KeywordTok{nullptr}\OperatorTok{);}
\NormalTok{glCompileShader}\OperatorTok{(}\NormalTok{shader}\OperatorTok{);}
\NormalTok{glGetShaderiv}\OperatorTok{(}\NormalTok{shader}\OperatorTok{,}\NormalTok{ GL\_COMPILE\_STATUS}\OperatorTok{,} \OperatorTok{\&}\NormalTok{ok}\OperatorTok{);}
\end{Highlighting}
\end{Shaded}

A shader \textbf{compiles} separately, then vertex and fragment shaders
are \textbf{linked} into one shader program. The project checks
compile/link status and throws an error containing the GPU driver's log
on failure. This differs from compiling C++ files with \texttt{g++}:
GLSL is compiled by the graphics driver at runtime.

\subsection{The packed vertex layout}\label{the-packed-vertex-layout}

\begin{Shaded}
\begin{Highlighting}[]
\KeywordTok{struct}\NormalTok{ Vertex }\OperatorTok{\{}
\NormalTok{    Vector3D position}\OperatorTok{;}
\NormalTok{    Vector3D colour}\OperatorTok{;}
\OperatorTok{\};}
\end{Highlighting}
\end{Shaded}

\texttt{vao} (vertex array object) stores how vertex attributes are
interpreted. \texttt{vbo} (vertex buffer object) owns the GPU-side byte
buffer containing vertex data. The code declares the positions/colours
to the GPU:

\begin{Shaded}
\begin{Highlighting}[]
\NormalTok{glBindVertexArray}\OperatorTok{(}\NormalTok{vao}\OperatorTok{);}
\NormalTok{glBindBuffer}\OperatorTok{(}\NormalTok{GL\_ARRAY\_BUFFER}\OperatorTok{,}\NormalTok{ vbo}\OperatorTok{);}
\NormalTok{glVertexAttribPointer}\OperatorTok{(}\DecValTok{0}\OperatorTok{,} \DecValTok{3}\OperatorTok{,}\NormalTok{ GL\_FLOAT}\OperatorTok{,}\NormalTok{ GL\_FALSE}\OperatorTok{,}
    \KeywordTok{sizeof}\OperatorTok{(}\NormalTok{Vertex}\OperatorTok{),} \KeywordTok{reinterpret\_cast}\OperatorTok{\textless{}}\DataTypeTok{void}\OperatorTok{*\textgreater{}(}\NormalTok{offsetof}\OperatorTok{(}\NormalTok{Vertex}\OperatorTok{,}\NormalTok{ position}\OperatorTok{)));}
\NormalTok{glEnableVertexAttribArray}\OperatorTok{(}\DecValTok{0}\OperatorTok{);}
\end{Highlighting}
\end{Shaded}

\begin{itemize}
\tightlist
\item
  First argument \texttt{0}: GLSL attribute location 0.
\item
  \texttt{3}: three scalar components per vertex.
\item
  \texttt{GL\_FLOAT}: underlying scalar data is \texttt{float}.
\item
  \texttt{GL\_FALSE}: don't normalise float values.
\item
  \texttt{sizeof(Vertex)}: stride, in bytes, between consecutive
  vertices.
\item
  \texttt{offsetof(...)}: byte offset from start of structure to member.
\item
  \texttt{reinterpret\_cast\textless{}void*\textgreater{}}: expresses
  the byte offset in the API's expected pointer-typed argument.
\end{itemize}

\textbf{Portability caveat:} for robust production code, ensure the
\texttt{Vertex} layout matches the buffer attribute definition,
including GLM alignment configuration. A future mesh format should
consider explicit tightly packed fields or static layout assertions.

\subsection{Per-frame rendering}\label{per-frame-rendering}

The renderer assembles a dynamic
\texttt{std::vector\textless{}Vertex\textgreater{}} for ground lines,
active beams, broken beams, mouse spring and points. For each group, the
local \texttt{drawVertices} lambda uploads current vertices:

\begin{Shaded}
\begin{Highlighting}[]
\NormalTok{glBufferData}\OperatorTok{(}\NormalTok{GL\_ARRAY\_BUFFER}\OperatorTok{,}
    \KeywordTok{static\_cast}\OperatorTok{\textless{}}\NormalTok{GLsizeiptr}\OperatorTok{\textgreater{}(}\NormalTok{vertices}\OperatorTok{.}\NormalTok{size}\OperatorTok{()} \OperatorTok{*} \KeywordTok{sizeof}\OperatorTok{(}\NormalTok{Vertex}\OperatorTok{)),}
\NormalTok{    vertices}\OperatorTok{.}\NormalTok{data}\OperatorTok{(),}\NormalTok{ GL\_DYNAMIC\_DRAW}\OperatorTok{);}
\NormalTok{glDrawArrays}\OperatorTok{(}\NormalTok{mode}\OperatorTok{,} \DecValTok{0}\OperatorTok{,} \KeywordTok{static\_cast}\OperatorTok{\textless{}}\NormalTok{GLsizei}\OperatorTok{\textgreater{}(}\NormalTok{vertices}\OperatorTok{.}\NormalTok{size}\OperatorTok{()));}
\end{Highlighting}
\end{Shaded}

\texttt{vertices.data()} is a pointer to contiguous C++ vector memory.
OpenGL copies data into the currently bound VBO. The
\texttt{GL\_DYNAMIC\_DRAW} usage hint indicates the data is expected to
change frequently. This is functional for 28 beams; at larger scales,
batch allocations and persistent buffers may become more efficient.

\subsection{Why OpenGL cleanup
matters}\label{why-opengl-cleanup-matters}

The renderer's destructor calls \texttt{glDeleteBuffers},
\texttt{glDeleteVertexArrays}, \texttt{glDeleteProgram}, destroys the
GLFW window, then terminates GLFW. It makes the context current before
deleting GL resources. This is \textbf{RAII} in practice: graphics
resources have a controlled lifetime, instead of relying on a Python GC
finalizer.

\subsection{Render order and depth}\label{render-order-and-depth}

The ground grid is drawn with depth testing. Beams are depth-tested as
lines. The source temporarily disables depth testing when drawing node
points so they remain easy to locate. That \textbf{intentionally lets
rear nodes remain visible through front geometry} for interaction and
can lead to ambiguous selections: a useful debug feature but not a
production occlusion policy.

\subsection{Future graphics
improvements}\label{future-graphics-improvements}

Static grid VBO, more deliberate GPU buffer allocation, solid mesh
skinning over the node cage, picking with depth, scalable vertex
colours/strain gradients, cameras, optional ImGui inspector, and
consistent frame pacing.

Sources:
\href{https://learnopengl.com/Getting-started/Hello-Triangle}{LearnOpenGL
Hello Triangle},
\href{https://learnopengl.com/Getting-started/Camera}{LearnOpenGL
Camera}, \href{https://www.glfw.org/docs/latest/quick_guide.html}{GLFW
quick guide}.

\newpage

\section{The GLEW Failure and GLAD 2
Migration}\label{the-glew-failure-and-glad-2-migration}

\subsection{What actually happened on 31 October
2026}\label{what-actually-happened-on-31-october-2026}

The user installed/upgraded packages using \texttt{pacman}. The first 3D
renderer still used GLEW. A build of the existing sources succeeded, but
running failed with:

\begin{Shaded}
\begin{Highlighting}[]
\NormalTok{SoloPhysics error: Failed to initialise GLEW: Missing GL version}
\end{Highlighting}
\end{Shaded}

A key debugging detail: a GLAD ZIP was referenced in instructions but
\textbf{not actually present in \texttt{\textasciitilde{}/Downloads}}.
The first \texttt{unzip} command failed, so running CMake in a folder
named \texttt{build-glad} still built the GLEW code
(\texttt{-\/-\ Found\ GLEW\ ...}). This was a \emph{patch-not-applied
problem}, independent of the underlying graphics loader error.

After the ZIP was downloaded/extracted and the project rebuilt, the user
confirmed: \textbf{``it works now''}.

\subsection{GLEW vs GLAD}\label{glew-vs-glad}

Both are OpenGL function loaders. Modern OpenGL entry points are
commonly retrieved at runtime from the active context, rather than being
assumed statically exported by a system header/library. GLFW creates the
window/context and provides \texttt{glfwGetProcAddress}. GLAD generates
bindings for the requested API version. In this project, GLAD was
selected to avoid the earlier GLEW initialization problem in the user's
GLFW/Wayland setup.

\begin{quote}
\textbf{Note:} Root-cause precision The error alone does \textbf{not}
prove Wayland/GLX was the sole technical cause. There was a suspected
mismatch between GLEW's initialization path and the native context on
that machine. The concrete verified facts are that GLEW failed, the old
code persisted while the patch was absent, and the GLAD build worked
after the patch was applied.
\end{quote}

\subsection{GLAD 2, not GLAD 1}\label{glad-2-not-glad-1}

The v0.4 source uses:

\begin{Shaded}
\begin{Highlighting}[]
\PreprocessorTok{\#include }\ImportTok{\textless{}glad/gl.h\textgreater{}}
\PreprocessorTok{\#include }\ImportTok{\textless{}GLFW/glfw3.h\textgreater{}}

\CommentTok{// Only after glfwMakeContextCurrent(window):}
\AttributeTok{const} \DataTypeTok{int}\NormalTok{ loadedVersion }\OperatorTok{=}\NormalTok{ gladLoadGL}\OperatorTok{(}\NormalTok{glfwGetProcAddress}\OperatorTok{);}
\ControlFlowTok{if} \OperatorTok{(}\NormalTok{loadedVersion }\OperatorTok{==} \DecValTok{0} \OperatorTok{||} \OperatorTok{!}\NormalTok{GLAD\_GL\_VERSION\_3\_3}\OperatorTok{)}
    \ControlFlowTok{throw} \BuiltInTok{std::}\NormalTok{runtime\_error}\OperatorTok{(}\StringTok{"Failed to load OpenGL 3.3 functions with GLAD"}\OperatorTok{);}
\end{Highlighting}
\end{Shaded}

\texttt{\textless{}glad/gl.h\textgreater{}} and \texttt{gladLoadGL(...)}
are the GLAD 2 API used here; don't substitute
\texttt{\textless{}glad/glad.h\textgreater{}} and
\texttt{gladLoadGLLoader(...)} examples from GLAD 1.

\textbf{Correct order}: initialize GLFW → choose 3.3 core hints → create
window → make context current → load GLAD → compile shader → create
VAO/VBO → draw. The context must exist and be current when GLAD obtains
entry points.

\subsection{Current CMake method}\label{current-cmake-method}

The file in the supplied ZIP pins an upstream GLAD Git commit and uses
\texttt{FetchContent}:

\begin{Shaded}
\begin{Highlighting}[]
\KeywordTok{include}\NormalTok{(}\ImportTok{FetchContent}\NormalTok{)}
\FunctionTok{FetchContent\_Declare}\NormalTok{(}
\NormalTok{    glad}
    \OtherTok{GIT\_REPOSITORY}\NormalTok{ https://github.com/Dav1dde/glad.git}
    \OtherTok{GIT\_TAG}\NormalTok{ 658f48e72aee3c6582e80b05ac0f8787a64fe6bb}
    \OtherTok{SOURCE\_SUBDIR}\NormalTok{ cmake}
    \OtherTok{GIT\_SHALLOW} \OtherTok{FALSE}
\NormalTok{)}
\FunctionTok{FetchContent\_MakeAvailable}\NormalTok{(glad)}
\FunctionTok{glad\_add\_library}\NormalTok{(glad\_gl\_core\_33 STATIC REPRODUCIBLE API gl:core=}\DataTypeTok{3.3}\NormalTok{ EXTENSIONS NONE)}
\end{Highlighting}
\end{Shaded}

This runs a generator (requires Python/Jinja2) and creates a CMake
library target \texttt{glad\_gl\_core\_33}. The pinned revision makes
source selection repeatable; generated code and toolchains can still
vary across environments. It downloads from GitHub on first configure
unless cached.

\subsection{Debug checklist}\label{debug-checklist}

\begin{Shaded}
\begin{Highlighting}[]
\CommentTok{\# Did the source actually change?}
\FunctionTok{grep} \AttributeTok{{-}nE} \StringTok{\textquotesingle{}GLEW|glad|glfw\textquotesingle{}}\NormalTok{ CMakeLists.txt}
\CommentTok{\# Is the new include present?}
\FunctionTok{grep} \AttributeTok{{-}nE} \StringTok{\textquotesingle{}glew.h|glad/gl.h|gladLoadGL\textquotesingle{}}\NormalTok{ src/rendering/Renderer.cpp}
\CommentTok{\# Configure cleanly in a NEW directory}
\FunctionTok{cmake} \AttributeTok{{-}S}\NormalTok{ . }\AttributeTok{{-}B}\NormalTok{ build{-}glad }\AttributeTok{{-}G}\NormalTok{ Ninja}
\FunctionTok{cmake} \AttributeTok{{-}{-}build}\NormalTok{ build{-}glad}
\ExtensionTok{./build{-}glad/solo\_physics}
\end{Highlighting}
\end{Shaded}

This case is a good lesson: \textbf{a new build directory does not
change the source code or dependencies by itself}. Confirm patch
extraction before debugging runtime loader behaviour.

References: \href{https://github.com/Dav1dde/glad}{GLAD upstream
repository},
\href{https://github.com/Dav1dde/glad/blob/glad2/cmake/GladConfig.cmake}{GLAD
CMake API},
\href{https://www.glfw.org/docs/latest/quick_guide.html}{GLFW guide},
\href{https://cmake.org/cmake/help/latest/module/FetchContent.html}{CMake
FetchContent}.

\newpage

\section{Line-by-Line CMake and Linking
Guide}\label{line-by-line-cmake-and-linking-guide}

\subsection{\texorpdfstring{The actual supplied
\texttt{CMakeLists.txt}}{The actual supplied CMakeLists.txt}}\label{the-actual-supplied-cmakelists.txt}

The current reference file can be opened in
\texttt{Source\ Snapshots/v0.4-GLAD/CMakeLists.txt}. Its key lines:

\begin{Shaded}
\begin{Highlighting}[]
\KeywordTok{cmake\_minimum\_required}\NormalTok{(}\OtherTok{VERSION} \DataTypeTok{3.23}\NormalTok{)}
\KeywordTok{project}\NormalTok{(SoloPhysics }\OtherTok{LANGUAGES} \OtherTok{C} \OtherTok{CXX}\NormalTok{)}
\KeywordTok{set}\NormalTok{(}\DecValTok{CMAKE\_CXX\_STANDARD}\NormalTok{ 20)}
\KeywordTok{set}\NormalTok{(}\DecValTok{CMAKE\_CXX\_STANDARD\_REQUIRED} \OtherTok{ON}\NormalTok{)}
\KeywordTok{set}\NormalTok{(}\DecValTok{CMAKE\_CXX\_EXTENSIONS} \OtherTok{OFF}\NormalTok{)}
\KeywordTok{find\_package}\NormalTok{(glm }\OtherTok{CONFIG} \OtherTok{REQUIRED}\NormalTok{)}
\KeywordTok{find\_package}\NormalTok{(glfw3 }\OtherTok{CONFIG} \OtherTok{REQUIRED}\NormalTok{)}
\KeywordTok{include}\NormalTok{(}\ImportTok{FetchContent}\NormalTok{)}
\FunctionTok{FetchContent\_Declare}\NormalTok{(glad}
    \OtherTok{GIT\_REPOSITORY}\NormalTok{ https://github.com/Dav1dde/glad.git}
    \OtherTok{GIT\_TAG}\NormalTok{ 658f48e72aee3c6582e80b05ac0f8787a64fe6bb}
    \OtherTok{SOURCE\_SUBDIR}\NormalTok{ cmake}
    \OtherTok{GIT\_SHALLOW} \OtherTok{FALSE}
\NormalTok{)}
\FunctionTok{FetchContent\_MakeAvailable}\NormalTok{(glad)}
\FunctionTok{glad\_add\_library}\NormalTok{(glad\_gl\_core\_33 STATIC REPRODUCIBLE API gl:core=}\DataTypeTok{3.3}\NormalTok{ EXTENSIONS NONE)}
\KeywordTok{add\_executable}\NormalTok{(}\BaseNTok{solo\_physics}
\NormalTok{    src/main.cpp}
\NormalTok{    src/physics/PhysicsWorld.cpp}
\NormalTok{    src/rendering/Renderer.cpp)}
\KeywordTok{target\_include\_directories}\NormalTok{(}\BaseNTok{solo\_physics} \OtherTok{PRIVATE} \DecValTok{$\{CMAKE\_CURRENT\_SOURCE\_DIR\}}\NormalTok{/include)}
\KeywordTok{target\_compile\_definitions}\NormalTok{(}\BaseNTok{solo\_physics} \OtherTok{PRIVATE}\NormalTok{ GLFW\_INCLUDE\_NONE)}
\KeywordTok{target\_link\_libraries}\NormalTok{(}\BaseNTok{solo\_physics} \OtherTok{PRIVATE} \BaseNTok{glm::glm}\NormalTok{ glfw glad\_gl\_core\_33)}
\KeywordTok{target\_compile\_options}\NormalTok{(}\BaseNTok{solo\_physics} \OtherTok{PRIVATE}\NormalTok{ {-}Wall {-}Wextra {-}Wpedantic)}
\end{Highlighting}
\end{Shaded}

\subsection{Decode it one line at a
time}\label{decode-it-one-line-at-a-time}

\begin{longtable}[]{@{}
  >{\raggedright\arraybackslash}p{(\columnwidth - 4\tabcolsep) * \real{0.3333}}
  >{\raggedright\arraybackslash}p{(\columnwidth - 4\tabcolsep) * \real{0.3333}}
  >{\raggedright\arraybackslash}p{(\columnwidth - 4\tabcolsep) * \real{0.3333}}@{}}
\toprule\noalign{}
\begin{minipage}[b]{\linewidth}\raggedright
Command
\end{minipage} & \begin{minipage}[b]{\linewidth}\raggedright
What it does
\end{minipage} & \begin{minipage}[b]{\linewidth}\raggedright
Python analogy
\end{minipage} \\
\midrule\noalign{}
\endhead
\bottomrule\noalign{}
\endlastfoot
\texttt{cmake\_minimum\_required(...)} & minimum CMake needed to
understand this build & minimum tool version \\
\texttt{project(...\ LANGUAGES\ C\ CXX)} & enable C and C++ compilation;
GLAD generates C source & packaging metadata + compilers \\
\texttt{CMAKE\_CXX\_STANDARD\ 20} & request C++20 & choosing
syntax/runtime target \\
\texttt{CMAKE\_CXX\_EXTENSIONS\ OFF} & prefer standard C++ over
GNU-specific C++ language mode & strict portability configuration \\
\texttt{find\_package(glm\ CONFIG\ REQUIRED)} & locate installed GLM
CMake package target & resolve external dependency \\
\texttt{find\_package(glfw3\ CONFIG\ REQUIRED)} & locate GLFW package
and target & resolve external dependency \\
\texttt{include(FetchContent)} & load built-in CMake dependency helper &
import build script functionality \\
\texttt{FetchContent\_Declare} & describe upstream Git source and
revision & pin package source URL + commit \\
\texttt{FetchContent\_MakeAvailable} & fetch/add dependency to configure
& obtain/install build dependency \\
\texttt{glad\_add\_library(...)} & generate GLAD bindings/library &
generate compile-time bindings \\
\texttt{add\_executable} & declare target and source files & define
application build \\
\texttt{target\_include\_directories} & where \texttt{\#include} looks
for local headers & interpreter import search path analogy \\
\texttt{target\_compile\_definitions} & define preprocessor flag &
build-time configuration flag \\
\texttt{target\_link\_libraries} & link the app to required libraries &
native linking, not \texttt{import} \\
\texttt{target\_compile\_options} & enable useful warnings &
lint/checker options \\
\end{longtable}

\subsection{CMake language syntax is not
C++}\label{cmake-language-syntax-is-not-c}

\texttt{set(X\ 20)} has parentheses but is a CMake \emph{command}, not a
C++ function. Values are usually \textbf{strings/arguments} in the CMake
language; variable interpolation uses \texttt{\$\{VARIABLE\}}, not
Python f-strings. \texttt{PRIVATE}/\texttt{PUBLIC}/\texttt{INTERFACE}
determine whether usage requirements propagate to depending targets, not
the C++ access modifiers inside a class.

\subsection{Compile versus link
errors}\label{compile-versus-link-errors}

A missing \texttt{\#include} file is a
\textbf{preprocessing/compilation} problem. A function declared but
never compiled into any linked object can cause
\texttt{undefined\ reference} at \textbf{link} time. Missing
\texttt{glfw} or generated \texttt{glad\_gl\_core\_33} linkage is also a
linking/configuration issue. Failure after the binary launches and
prints a GLFW/GLEW/GLAD error is a \textbf{runtime} issue.

\subsection{Rebuild best practices}\label{rebuild-best-practices}

\begin{Shaded}
\begin{Highlighting}[]
\CommentTok{\# Normal incremental build}
\FunctionTok{cmake} \AttributeTok{{-}{-}build}\NormalTok{ build{-}glad}

\CommentTok{\# Reconfigure when CMakeLists.txt changes}
\FunctionTok{cmake} \AttributeTok{{-}S}\NormalTok{ . }\AttributeTok{{-}B}\NormalTok{ build{-}glad }\AttributeTok{{-}G}\NormalTok{ Ninja}

\CommentTok{\# Prefer fresh independent directory to troubleshoot cache issues}
\FunctionTok{cmake} \AttributeTok{{-}S}\NormalTok{ . }\AttributeTok{{-}B}\NormalTok{ build{-}test }\AttributeTok{{-}G}\NormalTok{ Ninja}
\FunctionTok{cmake} \AttributeTok{{-}{-}build}\NormalTok{ build{-}test}
\end{Highlighting}
\end{Shaded}

Do not indiscriminately delete a directory that might contain something
you placed there manually. Generated build directories can generally be
recreated; keep source code outside them.

See
\href{https://cmake.org/cmake/help/latest/module/FetchContent.html}{official
CMake FetchContent documentation},
\href{https://github.com/Dav1dde/glad/blob/glad2/cmake/GladConfig.cmake}{GLAD
CMake definition}.

\newpage

\section{Real Errors, Diagnostics and
Lessons}\label{real-errors-diagnostics-and-lessons}

\subsection{Debugging timeline}\label{debugging-timeline}

\begin{longtable}[]{@{}
  >{\raggedright\arraybackslash}p{(\columnwidth - 6\tabcolsep) * \real{0.2500}}
  >{\raggedright\arraybackslash}p{(\columnwidth - 6\tabcolsep) * \real{0.2500}}
  >{\raggedright\arraybackslash}p{(\columnwidth - 6\tabcolsep) * \real{0.2500}}
  >{\raggedright\arraybackslash}p{(\columnwidth - 6\tabcolsep) * \real{0.2500}}@{}}
\toprule\noalign{}
\begin{minipage}[b]{\linewidth}\raggedright
Symptom
\end{minipage} & \begin{minipage}[b]{\linewidth}\raggedright
Observed cause / diagnosis
\end{minipage} & \begin{minipage}[b]{\linewidth}\raggedright
Fix
\end{minipage} & \begin{minipage}[b]{\linewidth}\raggedright
General lesson
\end{minipage} \\
\midrule\noalign{}
\endhead
\bottomrule\noalign{}
\endlastfoot
Changing damping seemed ineffective & A local \texttt{main.cpp} edit
wasn't saved before build & Save, then rebuild; print config values &
Verify inputs and build freshness \\
Bottom nodes sat at ground with positive vertical velocity & Repeated
tiny penetration plus restitution & Small-speed resting threshold &
Contact impulses require persistent resting treatment \\
Clicking seemed to do nothing & Small picking radius / unclear
mouse-to-world state & Print cursor world coords; increase radius; show
selected node & Inspect each link in input→selection→force chain \\
Dragging not acting through physics & Node force assigned before
\texttt{resetForces()} & Add callback inside step between reset and
integration & Put state updates in correct stage \\
Beams never turned orange/broke &
\texttt{plasticDeformationEnabled=false}, \texttt{breakStrain=100}
remained set & Enable plasticity, lower test thresholds, anchor a node &
Verify actual configuration, not the intended one \\
\texttt{i\ was\ not\ declared\ in\ this\ scope} &
\texttt{world.nodes{[}i{]}} written outside
\texttt{for(std::size\_t\ i...)} & Move fixed-node \texttt{continue}
inside indexed loop & Scope is lexical; \texttt{continue} is
loop-only \\
Copying instructions into \texttt{Renderer.cpp} caused invalid C++ &
Markdown explanatory text/code fences were pasted into \texttt{.cpp} &
Replace file with clean source & Code blocks ≠ prose instructions \\
\texttt{Failed\ to\ initialise\ GLEW:\ Missing\ GL\ version} & Runtime
loader/context incompatibility suspected & Migrate loader to GLAD 2 &
Distinguish loader from GPU/shader/physics \\
CMake still said \texttt{Found\ GLEW} after naming build folder
\texttt{build-glad} & Patch ZIP absent; \texttt{unzip} failed; source
still used GLEW & Actually download/extract patch, verify \texttt{grep}
results & Directory names do not change source \\
\end{longtable}

\subsection{Reproducible debugging
funnel}\label{reproducible-debugging-funnel}

\begin{Shaded}
\begin{Highlighting}[]
\NormalTok{flowchart TD}
\NormalTok{ A[Does it compile?] {-}{-}\textgreater{}|No| B[Read first compiler error]}
\NormalTok{ A {-}{-}\textgreater{}|Yes| C[Does it link?]}
\NormalTok{ C {-}{-}\textgreater{}|No| D[Check missing targets/source definitions]}
\NormalTok{ C {-}{-}\textgreater{}|Yes| E[Does window/context initialise?]}
\NormalTok{ E {-}{-}\textgreater{}|No| F[Check GLFW / loader / GPU versions]}
\NormalTok{ E {-}{-}\textgreater{}|Yes| G[Are physics numbers valid?]}
\NormalTok{ G {-}{-}\textgreater{}|No| H[Check dt, mass, forces, contact, strain]}
\NormalTok{ G {-}{-}\textgreater{}|Yes| I[Do pixels match state?]}
\NormalTok{ I {-}{-}\textgreater{}|No| J[Check shaders, VAO/VBO, matrices, depth]}
\NormalTok{ I {-}{-}\textgreater{}|Yes| K[Test input and interactivity]}
\end{Highlighting}
\end{Shaded}

\subsection{\texorpdfstring{What \texttt{ninja:\ no\ work\ to\ do}
actually
means}{What ninja: no work to do actually means}}\label{what-ninja-no-work-to-do-actually-means}

It means Ninja believes its declared inputs/outputs are up-to-date; it
does \textbf{not} prove your program uses the source you \emph{intended}
to copy. Verify path, file contents, CMake configure output and the
executable invoked. Use
\texttt{cmake\ -\/-build\ build-glad\ -\/-verbose} when necessary to
inspect compiler/linker commands.

\subsection{Warnings versus errors}\label{warnings-versus-errors}

The v0.4 build log included
\texttt{warning:\ variable\ \textquotesingle{}steps\textquotesingle{}\ set\ but\ not\ used}.
Warnings typically do not stop compilation with the current options;
they should still be addressed because they can reveal dead code or
mistaken state handling. \texttt{-Wall\ -Wextra\ -Wpedantic} increases
diagnostic visibility.

\subsection{Useful SoloLinux commands}\label{useful-sololinux-commands}

\begin{Shaded}
\begin{Highlighting}[]
\BuiltInTok{pwd}
\FunctionTok{find}\NormalTok{ include src }\AttributeTok{{-}maxdepth}\NormalTok{ 3 }\AttributeTok{{-}type}\NormalTok{ f }\KeywordTok{|} \FunctionTok{sort}
\ExtensionTok{rg} \AttributeTok{{-}n} \StringTok{\textquotesingle{}GLEW|glad|externalForceCallback|updateDamage\textquotesingle{}}\NormalTok{ include src CMakeLists.txt}
\FunctionTok{cmake} \AttributeTok{{-}{-}build}\NormalTok{ build{-}glad }\AttributeTok{{-}{-}verbose}
\ExtensionTok{./build{-}glad/solo\_physics}
\FunctionTok{git}\NormalTok{ diff }\AttributeTok{{-}{-}stat}
\FunctionTok{git}\NormalTok{ diff }\AttributeTok{{-}{-}}\NormalTok{ src/main.cpp}
\end{Highlighting}
\end{Shaded}

\texttt{rg} is ripgrep; use regular \texttt{grep\ -R} if unavailable.
The \texttt{git\ diff} commands inspect uncommitted changes.

\subsection{The rule: measure before
changing}\label{the-rule-measure-before-changing}

When debugging a simulator, log meaningful invariants: node count, beam
count, node mass, dt, position, velocity, force, rest length, strain,
broken count and plastic-history count. \textbf{Change one variable at a
time}. A powerful diagnostic order is: numerical state → force signs →
picking → rendering. Never use just the visual cube as a correctness
oracle.

\newpage

\section{Tests, Learning Exercises and Worked
Answers}\label{tests-learning-exercises-and-worked-answers}

\subsection{A. Regression test matrix}\label{a.-regression-test-matrix}

\begin{longtable}[]{@{}
  >{\raggedright\arraybackslash}p{(\columnwidth - 4\tabcolsep) * \real{0.3333}}
  >{\raggedright\arraybackslash}p{(\columnwidth - 4\tabcolsep) * \real{0.3333}}
  >{\raggedright\arraybackslash}p{(\columnwidth - 4\tabcolsep) * \real{0.3333}}@{}}
\toprule\noalign{}
\begin{minipage}[b]{\linewidth}\raggedright
Test
\end{minipage} & \begin{minipage}[b]{\linewidth}\raggedright
Setup
\end{minipage} & \begin{minipage}[b]{\linewidth}\raggedright
Expected trend/criterion
\end{minipage} \\
\midrule\noalign{}
\endhead
\bottomrule\noalign{}
\endlastfoot
Gravity-only drop & unconnected 1 kg node, no ground impact &
\(v_y\approx -9.81t\), \(y\approx y_0-\frac12gt^2\); discretisation
causes small error \\
Zero-force isolation & one node, gravity zero & position and velocity
stay unchanged \\
Equal/opposite beam loads & 2 free nodes linked by one stretched spring
& internal forces sum to approximately zero \\
Hanging spring equilibrium & anchor 3 m, rest 1 m, \(m=1\), \(k=100\) &
settled height approximately \textbf{1.9019 m} \\
Damping comparison & compare \(c=5\) versus \(c=0.6\) & smaller damping
oscillates longer \\
Square pre-impact & 4 nodes, equal masses, initial square & common
downward acceleration before contact \\
Resting ground & square after impact and settling & bottom \(y=0\), tiny
\(v_y\) after correction \\
Plasticity disabled & pull beam above \texttt{yieldStrain} & rest length
stays unchanged, history flag false \\
Plasticity enabled & pull above yield without breaking & rest length
changes; history flag becomes true \\
Fracture & pull past \texttt{breakStrain} & \texttt{broken=true}, no
spring force in future steps \\
GLAD startup & OpenGL 3.3 driver/context & logs a loaded version; shader
program links \\
3D picking & orbit around cube and click a node & selected node index
remains valid \\
\end{longtable}

\subsection{B. Exercise 1 --- Translate Python to
C++}\label{b.-exercise-1-translate-python-to-c}

Python:

\begin{Shaded}
\begin{Highlighting}[]
\NormalTok{mass }\OperatorTok{=} \FloatTok{1.0}
\NormalTok{position }\OperatorTok{=}\NormalTok{ [}\FloatTok{0.0}\NormalTok{, }\FloatTok{2.0}\NormalTok{, }\FloatTok{0.0}\NormalTok{]}
\NormalTok{velocity }\OperatorTok{=}\NormalTok{ [}\FloatTok{0.0}\NormalTok{, }\FloatTok{0.0}\NormalTok{, }\FloatTok{0.0}\NormalTok{]}
\end{Highlighting}
\end{Shaded}

Exercise: create a \texttt{Node} in C++ and set the same state.
\textbf{Answer:}

\begin{Shaded}
\begin{Highlighting}[]
\NormalTok{Node node}\OperatorTok{;}
\NormalTok{node}\OperatorTok{.}\NormalTok{mass }\OperatorTok{=} \FloatTok{1.0}\BuiltInTok{f}\OperatorTok{;}
\NormalTok{node}\OperatorTok{.}\NormalTok{position }\OperatorTok{=}\NormalTok{ Vector3D}\OperatorTok{\{}\FloatTok{0.0}\BuiltInTok{f}\OperatorTok{,} \FloatTok{2.0}\BuiltInTok{f}\OperatorTok{,} \FloatTok{0.0}\BuiltInTok{f}\OperatorTok{\};}
\NormalTok{node}\OperatorTok{.}\NormalTok{velocity }\OperatorTok{=}\NormalTok{ Vector3D}\OperatorTok{\{}\FloatTok{0.0}\BuiltInTok{f}\OperatorTok{\};}
\end{Highlighting}
\end{Shaded}

\subsection{C. Exercise 2 --- Calculate spring
force}\label{c.-exercise-2-calculate-spring-force}

A 1 m rest beam is 1.02 m long, \(k=650\) N/m, axial relative velocity
0.1 m/s, \(c=12\) N·s/m. Find scalar force. \textbf{Answer:}

\[F=650(0.02)+12(0.1)=13+1.2=14.2\,\mathrm N.\]

This assumes \texttt{direction} points A→B and positive relative
velocity means endpoints are separating.

\subsection{D. Exercise 3 --- Interpret the
lambda}\label{d.-exercise-3-interpret-the-lambda}

\begin{Shaded}
\begin{Highlighting}[]
\KeywordTok{auto}\NormalTok{ fn }\OperatorTok{=} \OperatorTok{[\&](}\NormalTok{PhysicsWorld}\OperatorTok{\&}\NormalTok{ world}\OperatorTok{)} \OperatorTok{\{}
    \ControlFlowTok{if} \OperatorTok{(!}\NormalTok{grabbedNode}\OperatorTok{)} \ControlFlowTok{return}\OperatorTok{;}
\NormalTok{    world}\OperatorTok{.}\NormalTok{nodes}\OperatorTok{[*}\NormalTok{grabbedNode}\OperatorTok{].}\NormalTok{force }\OperatorTok{+=}\NormalTok{ Vector3D}\OperatorTok{\{}\DecValTok{1}\OperatorTok{,} \DecValTok{0}\OperatorTok{,} \DecValTok{0}\OperatorTok{\};}
\OperatorTok{\};}
\end{Highlighting}
\end{Shaded}

Explain why the compiler allows access to \texttt{grabbedNode} even
though it is not a parameter. \textbf{Answer:} \texttt{{[}\&{]}}
captures used enclosing local variables by reference; \texttt{fn} can
read the latest \texttt{grabbedNode}, assuming its lifetime remains
valid. \texttt{std::optional} presence is tested before dereferencing.

\subsection{E. Exercise 4 --- Derive static
equilibrium}\label{e.-exercise-4-derive-static-equilibrium}

Change the hanging-mass spring to \(k=200\) N/m, \(m=1\) kg, same
anchor/rest length. Predict position. \textbf{Answer:} extension
\(mg/k=0.04905\) m, so

\[y_{eq}=3-(1+0.04905)=1.95095\,\mathrm m.\]

\subsection{F. Exercise 5 --- Why 28
beams?}\label{f.-exercise-5-why-28-beams}

For eight distinct nodes, an edge between every unordered pair creates
\(8\times7/2=28\) beams. This is the complete graph \(K_8\). Not all 28
are necessary for a rigid braced cube; the current topology is a simple
heavily reinforced learning structure.

\subsection{G. Exercise 6 --- Debug a missing
force}\label{g.-exercise-6-debug-a-missing-force}

You add \texttt{world.nodes{[}3{]}.force\ +=\ F;} just before calling
\texttt{world.step(dt)}, but the node doesn't respond. Where does it go
wrong? \textbf{Answer:} \texttt{resetForces()} zeros it inside the
timestep; supply it via \texttt{externalForceCallback} (or a formal
persistent external-force input system) before beam
solution/integration.

\subsection{H. Exercise 7 --- Change OpenGL
safely}\label{h.-exercise-7-change-opengl-safely}

Explain why these lines must be ordered:

\begin{Shaded}
\begin{Highlighting}[]
\NormalTok{glfwMakeContextCurrent}\OperatorTok{(}\NormalTok{window}\OperatorTok{);}
\DataTypeTok{int}\NormalTok{ loadedVersion }\OperatorTok{=}\NormalTok{ gladLoadGL}\OperatorTok{(}\NormalTok{glfwGetProcAddress}\OperatorTok{);}
\end{Highlighting}
\end{Shaded}

\textbf{Answer:} OpenGL entry-point lookup depends on a current context.
Trying to load before binding a context may fail, even if the compiler
and linker are happy.

\subsection{I. Suggested exercise
roadmap}\label{i.-suggested-exercise-roadmap}

\begin{enumerate}
\def\labelenumi{\arabic{enumi}.}
\tightlist
\item
  Write a tiny C++ test executable that compiles only
  \texttt{PhysicsWorld.cpp} plus tests; no renderer.
\item
  Compare analytical and simulated oscillator periods over multiple
  timesteps.
\item
  Record total energy (kinetic + spring + gravitational) with no
  damping/ground and explain integration drift.
\item
  Visualise active-beam strain as a continuous gradient while keeping
  plastic history a distinct overlay.
\item
  Expose material parameters in a configuration file rather than
  hardcoding each scenario.
\item
  Add mouse-input replay so regression tests don't depend on how quickly
  you move the mouse.
\end{enumerate}

Each task depends on concepts already encountered in the project; keep a
lab note recording hypothesis, method, result and discrepancies.

\newpage

\section{Current Limits and a Path Toward Vehicle
Simulation}\label{current-limits-and-a-path-toward-vehicle-simulation}

\subsection{Present capabilities (confirmed by
source/conversation)}\label{present-capabilities-confirmed-by-sourceconversation}

\begin{itemize}
\tightlist
\item
  3D point-mass nodes and axial spring--damper beams.
\item
  External gravity; fixed anchor nodes.
\item
  Semi-implicit Euler, fixed-step application loop.
\item
  Simplified flat-ground node contact and restitution.
\item
  Yield-triggered rest-length change, broken-beam flags.
\item
  OpenGL 3.3 line/point rendering, orbit/pan/zoom, screen-space picking
  and mouse spring dragging.
\item
  Native GLFW context with GLAD 2; confirmed launched successfully after
  migration.
\end{itemize}

\subsection{Not yet established}\label{not-yet-established}

A working 3D cube display is \textbf{not} the same as validated
soft-body vehicle dynamics. Current limitations include: simplistic
contact, no self-collision, no continuous collision detection, no proper
constraints solver, no beam bending/buckling/torsion, no traction or
suspension/wheels, no vehicle drivetrain, no contact patches, no
material calibration, no deformable surface meshes, no broad-phase
spatial partitioning, no replay-determinism guarantee, and no numerical
performance benchmarks for thousands of beams.

\subsection{Suggested engineering
order}\label{suggested-engineering-order}

\begin{longtable}[]{@{}
  >{\raggedright\arraybackslash}p{(\columnwidth - 4\tabcolsep) * \real{0.3333}}
  >{\raggedright\arraybackslash}p{(\columnwidth - 4\tabcolsep) * \real{0.3333}}
  >{\raggedright\arraybackslash}p{(\columnwidth - 4\tabcolsep) * \real{0.3333}}@{}}
\toprule\noalign{}
\begin{minipage}[b]{\linewidth}\raggedright
Phase
\end{minipage} & \begin{minipage}[b]{\linewidth}\raggedright
Priority
\end{minipage} & \begin{minipage}[b]{\linewidth}\raggedright
Why
\end{minipage} \\
\midrule\noalign{}
\endhead
\bottomrule\noalign{}
\endlastfoot
1 & Unit/regression testing & Guard physics correctness while changing
architecture \\
2 & Separate simulation library, scenario tests, UI & Easier to
maintain/test than one large \texttt{main.cpp} \\
3 & More robust ground contact and friction & Avoid velocity clamping
artefacts \\
4 & Better strain/damage history & Separate elastic/plastic/total
strain; test energy and break criteria \\
5 & Collision geometry and broad phase & Objects can actually crash into
each other \\
6 & Vehicle chassis + basic suspension & First vehicle-like dynamics \\
7 & Surface mesh attached to node cage & Visible deformation beyond line
segments \\
8 & Constraint solvers, possibly XPBD & Address stiffness/time-step
coupling and robust contacts \\
9 & Profiling and scalability & Support realistic node/beam counts \\
\end{longtable}

An eventual switch to \textbf{XPBD} or another constraint approach is a
research direction, not an assertion that the current spring formulation
is wrong. XPBD specifically targets problems of stiffness depending on
timestep/solver iteration count; see Macklin et al.,
\href{https://mmacklin.com/xpbd.pdf}{XPBD paper}.

\subsection{Research questions to answer
experimentally}\label{research-questions-to-answer-experimentally}

\begin{enumerate}
\def\labelenumi{\arabic{enumi}.}
\tightlist
\item
  How large can \texttt{stiffness} become at 240 Hz before noticeable
  instability?
\item
  How does the restitution threshold change the energy budget?
\item
  Is fracture triggered by instantaneous strain, accumulated history,
  stress, or energy release?
\item
  Does the rendering pick the nearest \emph{visible} node, or just
  nearest in screen-space among all nodes?
\item
  How many nodes and beams can be simulated at 240 Hz on the target
  hardware?
\item
  How will a vehicle's tyres/suspension interact with the same contact
  system?
\end{enumerate}

\subsection{Source discipline}\label{source-discipline}

Official BeamNG descriptions are a \textbf{conceptual inspiration} for
node--beam simulation, but the SoloPhysics code remains its own
simplified educational implementation. Do not infer BeamNG internals
from matching terms or copy parameter units without verification.

\newpage

\section{Websites, Sources and
Attribution}\label{websites-sources-and-attribution}

\subsection{Provenance statement ---
important}\label{provenance-statement-important}

This bibliography has \textbf{two distinct classes of evidence}:

\begin{itemize}
\tightlist
\item
  \textbf{Project/author evidence (P):} user-supplied code, build logs,
  outputs, and conversation milestones. Those sources establish what
  SoloPhysics \emph{actually did}.
\item
  \textbf{External learning/reference (R):} authoritative websites
  checked while preparing this documentation (31 Oct 2026) or explicitly
  recommended/discussed during development. These explain third-party
  APIs and algorithms but are \textbf{not proof} that every page was
  personally consulted at the moment the original code was authored. The
  earlier assistant did not preserve an exhaustive browser-history
  ledger; therefore \textbf{an exhaustive list of every website
  previously visited cannot honestly be certified}.
\end{itemize}

This distinction avoids retroactively claiming websites were used when
the evidence only shows they are relevant/recommended.

\subsection{Direct project evidence}\label{direct-project-evidence}

\begin{longtable}[]{@{}
  >{\raggedright\arraybackslash}p{(\columnwidth - 4\tabcolsep) * \real{0.3333}}
  >{\raggedright\arraybackslash}p{(\columnwidth - 4\tabcolsep) * \real{0.3333}}
  >{\raggedright\arraybackslash}p{(\columnwidth - 4\tabcolsep) * \real{0.3333}}@{}}
\toprule\noalign{}
\begin{minipage}[b]{\linewidth}\raggedright
ID
\end{minipage} & \begin{minipage}[b]{\linewidth}\raggedright
Evidence
\end{minipage} & \begin{minipage}[b]{\linewidth}\raggedright
Use
\end{minipage} \\
\midrule\noalign{}
\endhead
\bottomrule\noalign{}
\endlastfoot
P1 & Conversations and recorded milestone notes, 28--31 Oct 2026 &
Design decisions, progression, suspended-mass and square-drop
measurements \\
P2 & User's pasted \texttt{main.cpp} and \texttt{Renderer.cpp} excerpts
from 2D iterations & Actual C++ syntax, callbacks and earlier
debugging \\
P3 & CLI outputs for hanging mass, square drop, and improved contact &
Numerical results (\(1.9019\) m and resting square) \\
P4 & User's build logs, 31 Oct 2026: missing patch ZIP, GLEW failure &
Loader transition diagnosis \\
P5 & \texttt{SoloPhysics-v0.4-GLAD-upgrade.zip} bundle available in the
conversation & Verified v0.4 \texttt{CMakeLists.txt}, \texttt{main.cpp},
\texttt{Renderer.hpp}, \texttt{Renderer.cpp} \\
P6 & User statement, 31 Oct 2026, ``it works now'' & GLAD renderer runs
on user's machine \\
\end{longtable}

\textbf{Source completeness limit:} the current live repository and its
newest local \texttt{PhysicsWorld.cpp} were not read in their entirety
for this handbook. Source snapshots are the supplied GLAD upgrade
bundle, not a Git clone of the user's project.

\subsection{Official documentation and educational
sources}\label{official-documentation-and-educational-sources}

\begin{longtable}[]{@{}
  >{\raggedright\arraybackslash}p{(\columnwidth - 6\tabcolsep) * \real{0.2500}}
  >{\raggedright\arraybackslash}p{(\columnwidth - 6\tabcolsep) * \real{0.2500}}
  >{\raggedright\arraybackslash}p{(\columnwidth - 6\tabcolsep) * \real{0.2500}}
  >{\raggedright\arraybackslash}p{(\columnwidth - 6\tabcolsep) * \real{0.2500}}@{}}
\toprule\noalign{}
\begin{minipage}[b]{\linewidth}\raggedright
ID
\end{minipage} & \begin{minipage}[b]{\linewidth}\raggedright
Website / precise URL
\end{minipage} & \begin{minipage}[b]{\linewidth}\raggedright
Why it matters
\end{minipage} & \begin{minipage}[b]{\linewidth}\raggedright
Status
\end{minipage} \\
\midrule\noalign{}
\endhead
\bottomrule\noalign{}
\endlastfoot
R1 & \href{https://www.beamng.com/game/about/physics/}{BeamNG: Soft-body
Physics} & Nodes, massless beams, axial constraints,
deformation/breaking & verified reference; conceptual inspiration \\
R2 & \href{https://docs.beamng.com/modding/vehicle/intro_jbeam/}{BeamNG:
Introduction to JBeam} & Structure, node/beam modelling, orientation &
verified reference \\
R3 & \href{https://www.glfw.org/docs/latest/quick_guide.html}{GLFW
Getting Started} & Window/context creation, rendering loop, GLAD loader
ordering & verified reference \\
R4 & \href{https://www.glfw.org/docs/latest/input.html}{GLFW Input
Guide} & Mouse/key polling, callbacks and window user pointers &
verified reference \\
R5 & \href{https://github.com/Dav1dde/glad}{GLAD upstream repository} &
Generator and GLAD 2 function loading & verified reference; exact
repository pinned by project \\
R6 &
\href{https://github.com/Dav1dde/glad/blob/glad2/cmake/GladConfig.cmake}{GLAD
CMake library definition} &
\texttt{glad\_add\_library(...\ API\ gl:core=3.3)} & verified
reference \\
R7 &
\href{https://cmake.org/cmake/help/latest/module/FetchContent.html}{CMake
FetchContent} & Fetch/build upstream dependency & verified reference \\
R8 & \href{https://github.com/g-truc/glm/blob/master/manual.md}{GLM
Manual} & \texttt{glm::vec3}, \texttt{glm::lookAt},
\texttt{glm::perspective}, \texttt{glm::project/unProject} & verified
reference \\
R9 &
\href{https://learnopengl.com/Getting-started/Hello-Triangle}{LearnOpenGL:
Hello Triangle} & VBO, VAO, shaders and draw calls & verified
educational reading \\
R10 & \href{https://learnopengl.com/Getting-started/Camera}{LearnOpenGL:
Camera} & View matrices and camera motion & verified educational
reading \\
R11 & \href{https://gafferongames.com/post/fix_your_timestep/}{Gaffer on
Games: Fix Your Timestep!} & Fixed-step accumulator rationale & verified
educational reading \\
R12 &
\href{https://en.cppreference.com/cpp/utility/optional}{cppreference:
std::optional} & Value-or-no-value state for selection/drag & verified
C++ reference \\
R13 &
\href{https://en.cppreference.com/cpp/utility/functional/function}{cppreference:
std::function} & Stored external force callback & verified C++
reference \\
R14 &
\href{https://www.learncpp.com/cpp-tutorial/introduction-to-lambdas-anonymous-functions/}{LearnCpp:
Introduction to lambdas} & Captures, anonymous callable objects &
verified learning reference \\
R15 & \href{https://mmacklin.com/xpbd.pdf}{Macklin, Müller and
Chentanez: XPBD (2016)} & Future constraint-solver research; not current
algorithm & verified primary paper \\
\end{longtable}

All external links were checked or retrieved during this documentation
pass on \textbf{31 October 2026}. The exact content of external pages
may change later. For authoritative specifications, prefer upstream
docs; for tutorials, compare explanations against upstream APIs and
compiler output.

\subsection{Licence/attribution
awareness}\label{licenceattribution-awareness}

SoloPhysics uses third-party libraries. GLM, GLFW, GLAD, OpenGL and
CMake are distinct projects with their own documentation, licensing and
APIs. This handbook paraphrases their concepts and contains small
original instructional examples and user-project excerpts. It does
\textbf{not} reproduce substantial portions of external tutorials or
articles.

\subsection{Sources mentioned in original planning but not
implemented}\label{sources-mentioned-in-original-planning-but-not-implemented}

Early planning discussed position-based dynamics (PBD/XPBD) as a
possible future strategy. It is not the solver currently in SoloPhysics.
BeamNG references describe conceptual inspiration but the project has no
claim of implementing the complete BeamNG algorithm.

\newpage

\section{C++ Syntax Cheat Sheet for Python
Programmers}\label{c-syntax-cheat-sheet-for-python-programmers}

\subsection{\texorpdfstring{Read this before opening
\texttt{Renderer.cpp}}{Read this before opening Renderer.cpp}}\label{read-this-before-opening-renderer.cpp}

\begin{longtable}[]{@{}
  >{\raggedright\arraybackslash}p{(\columnwidth - 4\tabcolsep) * \real{0.3333}}
  >{\raggedright\arraybackslash}p{(\columnwidth - 4\tabcolsep) * \real{0.3333}}
  >{\raggedright\arraybackslash}p{(\columnwidth - 4\tabcolsep) * \real{0.3333}}@{}}
\toprule\noalign{}
\begin{minipage}[b]{\linewidth}\raggedright
C++ notation
\end{minipage} & \begin{minipage}[b]{\linewidth}\raggedright
Plain English
\end{minipage} & \begin{minipage}[b]{\linewidth}\raggedright
Python rough analogy
\end{minipage} \\
\midrule\noalign{}
\endhead
\bottomrule\noalign{}
\endlastfoot
\texttt{\#include\ \textless{}file\textgreater{}} & make declarations
from header available & \texttt{import} (compile-time, not runtime) \\
\texttt{\#pragma\ once} & don't include same header twice in translation
unit & no direct equivalent \\
\texttt{namespace\ solophysics} & scope for project symbols &
module/namespace \\
\texttt{::} & qualify a scope & \texttt{module.name} \\
\texttt{.member} & member of a value/reference & \texttt{obj.attr} \\
\texttt{ptr-\textgreater{}member} & member through a pointer &
\texttt{obj.attr} after pointer dereference \\
\texttt{\&} after a type & reference / alias & implicit object
reference, not identical \\
\texttt{\&x} expression & take address of x & no everyday equivalent \\
\texttt{*ptr} expression & dereference pointer & no everyday
equivalent \\
\texttt{*opt} & retrieve contained optional value & unwrap
\texttt{Optional} \\
\texttt{const\ T\&} & read-only borrowed alias & immutable-access view
analogy \\
\texttt{T\&\&} & rvalue reference & advanced, not yet needed \\
\texttt{auto} & compiler infers a concrete type & inferred typing, but
static \\
\texttt{constexpr} & compile-time value when possible & \texttt{Final}
isn't exact equivalent \\
\texttt{std::vector\textless{}T\textgreater{}} & dynamically sized
contiguous sequence & \texttt{list{[}T{]}} \\
\texttt{std::array\textless{}T,N\textgreater{}} & fixed-size array &
fixed-length tuple analogy \\
\texttt{std::size\_t} & unsigned size/index type & \texttt{int} for
indices \\
\texttt{std::optional\textless{}T\textgreater{}} & maybe T; maybe no
value & \texttt{T\ \textbar{}\ None} \\
\texttt{std::function\textless{}R(Args...)\textgreater{}} & stored
callable with signature & \texttt{Callable{[}{[}Args{]},\ R{]}} \\
\texttt{{[}{]}(...)\ \{...\}} & lambda, no captures & \texttt{lambda}
(but C++ allows block body) \\
\texttt{{[}\&{]}(...)\ \{...\}} & lambda captures used locals by
reference & closure captures variable bindings \\
\texttt{{[}={]}(...)\ \{...\}} & lambda captures used locals by value &
closure capture snapshot analogy \\
\texttt{{[}this{]}} & capture current object pointer & bound method
analogy \\
\texttt{for\ (auto\&\ x\ :\ xs)} & iterate by reference &
\texttt{for\ x\ in\ xs:} \\
\texttt{for\ (i=0;\ i\textless{}n;\ ++i)} & index loop &
\texttt{for\ i\ in\ range(n):} \\
\texttt{if\ (!flag)} & if not true & \texttt{if\ not\ flag:} \\
\texttt{\&\&}, \texttt{\textbar{}\textbar{}} & logical AND, OR &
\texttt{and}, \texttt{or} \\
\texttt{true}, \texttt{false} & boolean values & \texttt{True},
\texttt{False} \\
\texttt{nullptr} & null pointer & \texttt{None} (not the same type) \\
\texttt{\{0,1,2\}} & aggregate/brace initialisation & list/tuple literal
analogy \\
\texttt{std::cout\ \textless{}\textless{}\ x} & write stdout &
\texttt{print(x)} \\
\texttt{std::cerr\ \textless{}\textless{}\ x} & write stderr &
\texttt{print(x,\ file=sys.stderr)} \\
\texttt{throw} / \texttt{catch} & exception handling & \texttt{raise} /
\texttt{except} \\
\texttt{static\_cast\textless{}T\textgreater{}(x)} & explicit conversion
& \texttt{T(x)} (not exact) \\
\texttt{sizeof(T)} & number of bytes for C++ object type &
\texttt{struct.calcsize} analogy \\
\texttt{offsetof(T,member)} & member byte offset & C ABI \texttt{ctypes}
analogy \\
\texttt{R"TAG(...)TAG"} & raw string literal & \texttt{r"..."} but with
delimiter \\
\texttt{\textasciitilde{}Renderer()} & destructor & \texttt{\_\_del\_\_}
analogy, but timing differs \\
\texttt{=\ delete} & disable a function/operation & no direct
equivalent \\
\end{longtable}

\subsection{Special vector syntax}\label{special-vector-syntax}

\begin{Shaded}
\begin{Highlighting}[]
\KeywordTok{using}\NormalTok{ Vector3D }\OperatorTok{=}\NormalTok{ glm}\OperatorTok{::}\NormalTok{vec3}\OperatorTok{;}
\NormalTok{Vector3D p}\OperatorTok{\{}\FloatTok{1.0}\BuiltInTok{f}\OperatorTok{,} \FloatTok{2.0}\BuiltInTok{f}\OperatorTok{,} \FloatTok{3.0}\BuiltInTok{f}\OperatorTok{\};}
\NormalTok{Vector3D v}\OperatorTok{\{}\FloatTok{0.0}\BuiltInTok{f}\OperatorTok{\};}
\NormalTok{Vector3D moved }\OperatorTok{=}\NormalTok{ p }\OperatorTok{+}\NormalTok{ v }\OperatorTok{*} \FloatTok{0.1}\BuiltInTok{f}\OperatorTok{;}
\DataTypeTok{float}\NormalTok{ distance }\OperatorTok{=}\NormalTok{ glm}\OperatorTok{::}\NormalTok{length}\OperatorTok{(}\NormalTok{moved }\OperatorTok{{-}}\NormalTok{ p}\OperatorTok{);}
\DataTypeTok{float}\NormalTok{ alignedSpeed }\OperatorTok{=}\NormalTok{ glm}\OperatorTok{::}\NormalTok{dot}\OperatorTok{(}\NormalTok{v}\OperatorTok{,}\NormalTok{ glm}\OperatorTok{::}\NormalTok{normalize}\OperatorTok{(}\NormalTok{p}\OperatorTok{));}
\end{Highlighting}
\end{Shaded}

These operations are built by GLM's overloaded operators, analogous to
NumPy array arithmetic---not automatic features of built-in C arrays.

\subsection{Common traps}\label{common-traps}

\begin{itemize}
\tightlist
\item
  \texttt{=} assigns; \texttt{==} compares.
\item
  \texttt{\&} is reference in a declaration but address-of in many
  expressions.
\item
  \texttt{*} multiplies numbers and dereferences pointers/optional
  values depending on context.
\item
  \texttt{{[}{]}} is indexing for a container and capture syntax for a
  lambda.
\item
  \texttt{std::vector} references can be invalidated by reallocation.
\item
  \texttt{float} and \texttt{double} are not Python's single ordinary
  float type; literals ending in \texttt{f} are \texttt{float}.
\item
  Always check the \emph{first} compiler error before the cascade that
  follows.
\end{itemize}

\subsection{Further reading}\label{further-reading}

03 - C++ for a Python Developer • 04 - Data Model and Source
Architecture • 13 - Modern OpenGL and GLSL •
\href{https://www.learncpp.com/}{LearnCpp} •
\href{https://en.cppreference.com/}{cppreference}

\newpage

\end{document}
